\chapter{狭义相对论力学}

\section{洛伦兹变换}

\begin{problem}{12.1.1}
简单叙述一下需要狭义相对论来解决的难题;叙述一种可能不需要狭义相对论的早期理论,并举出一个证明这种理论是错误的实验;叙述一个证明狭义相对论可信的近代实验.
\end{problem}
\clearpage


\begin{problem}{12.1.2}
菲佐实验所用的装置如图 12.1 所示. 光源发出的、在真空中波长为 $\lambda$ 的光, 经镀有半透膜的玻璃片 M 分成两束相干光, 一束透射, 经 $\mathbf{M}_1, \mathbf{M}_2, \mathbf{M}_3$ , 再在 M 处透射, 射入望远镜 T; 另一束经 M 反射, 经 $\mathbf{M}_3, \mathbf{M}_2, \mathbf{M}_1$ , 再在 M 处反射, 射入望远镜 T, 前者顺水速而行, 后者逆水速而行. 设水速为 $v$ , 水的折射率为 $n$ , 若光相对于静止水中以太的速率为 $\frac{c}{n}$ , 以 $v$ 的速率运动的水带动以太的速率为 $kv$ , 用钠黄光作光源, $\lambda = 5893 \times 10^{-10} \mathrm{~m}, l = 1.5 \mathrm{~m}, v = 7.0 \mathrm{~m/s}, n = 1.33$ , 观察到与 $v = 0$ 的情况相比, 测得移动的干涉条纹为 $\delta = 0.19$ , 根据以太存在的假设, $k$ 应取何值?
\begin{center}
\begin{tikzpicture}[>=Stealth, line width=0.7pt]
  % Water tube assembly (U-shaped channel)
  \draw (2.40,1.84) -- (3.64,1.84) -- (3.64,3.26) -- (2.40,3.26);
  \draw (4.00,1.84) -- (3.80,1.84) -- (3.80,3.26) -- (4.00,3.26);
  \draw (2.40,1.16) -- (2.62,1.16) -- (2.62,0.90);
  \draw (2.76,0.90) -- (2.76,1.16) -- (4.00,1.16);
  \draw (2.40,3.94) -- (2.62,3.94) -- (2.62,4.20);
  \draw (2.76,4.20) -- (2.76,3.94) -- (4.00,3.94);

  % Water stream dashed lines and flow velocity arrows
  \draw[dash pattern=on 4pt off 3pt] (2.52,1.76) -- (3.05,1.76);
  \draw[dash pattern=on 4pt off 3pt] (3.42,1.76) -- (3.60,1.76);
  \draw[dash pattern=on 4pt off 3pt] (2.52,1.28) -- (3.88,1.28);
  \draw[dash pattern=on 4pt off 3pt] (2.52,3.82) -- (3.88,3.82);
  \draw[dash pattern=on 4pt off 3pt] (2.52,3.34) -- (3.05,3.34);
  \draw[dash pattern=on 4pt off 3pt] (3.42,3.34) -- (3.60,3.34);
  \draw[->] (2.98,1.59) -- (3.48,1.59);
  \node[above=1pt, inner sep=1pt] at (3.23,1.59) {$v$};
  \draw[->] (3.48,3.51) -- (2.98,3.51);
  \node[below=1pt, inner sep=1pt] at (3.23,3.51) {$v$};

  % Glass end-caps on tubes
  \draw[fill=white] (2.30,1.08) rectangle (2.40,1.92);
  \draw[fill=white] (4.00,1.08) rectangle (4.10,1.92);
  \draw[fill=white] (2.30,3.18) rectangle (2.40,4.02);
  \draw[fill=white] (4.00,3.18) rectangle (4.10,4.02);

  % Inlet and outlet arrows + labels
  \draw[->] (2.69,0.68) -- (2.69,1.10);
  \node[right=2pt] at (2.78,0.86) {进水口};
  \draw[->] (2.69,4.00) -- (2.69,4.42);
  \node[right=2pt] at (2.78,4.24) {出水口};

  % Length l dimension below lower tube
  \draw[thin] (2.40,0.52) -- (2.40,0.22);
  \draw[thin] (4.00,0.52) -- (4.00,0.22);
  \draw[<->] (2.40,0.37) -- (4.00,0.37) node[midway, fill=white, inner sep=3pt] {$l$};

  % Half-silvered mirror M (drawn before ray so ray passes cleanly through)
  \begin{scope}[shift={(1.30,1.50)}, rotate=45]
    \draw[fill=white] (-0.34,-0.045) rectangle (0.34,0.045);
  \end{scope}
  \node[above left=4pt] at (1.22,1.62) {$\mathrm{M}$};

  % Light ray paths
  \draw (0.05,1.50) -- (5.05,1.50) -- (5.05,3.60) -- (1.30,3.60) -- (1.30,0.36);

  % Light source S
  \foreach \ang in {0,45,90,135} {
    \draw (0.05,1.50) ++(\ang:0.12) -- ++(\ang+180:0.24);
  }
  \node[above=4pt] at (0.05,1.50) {$\mathrm{S}$};
  \node[below=5pt] at (0.05,1.50) {光源};

  % Mirrors M1, M2, M3 with back-side hatching ticks
  \begin{scope}[shift={(5.05,1.50)}, rotate=45]
    \draw[line width=0.85pt] (-0.34,0) -- (0.34,0);
    \foreach \x in {-0.26,-0.13,0,0.13,0.26} {
      \draw (\x,0) -- (\x-0.08,-0.10);
    }
  \end{scope}
  \node[below right=4pt] at (5.15,1.38) {$\mathrm{M}_1$};

  \begin{scope}[shift={(5.05,3.60)}, rotate=-45]
    \draw[line width=0.85pt] (-0.34,0) -- (0.34,0);
    \foreach \x in {-0.26,-0.13,0,0.13,0.26} {
      \draw (\x,0) -- (\x+0.08,0.10);
    }
  \end{scope}
  \node[above right=4pt] at (5.15,3.70) {$\mathrm{M}_2$};

  \begin{scope}[shift={(1.30,3.60)}, rotate=45]
    \draw[line width=0.85pt] (-0.34,0) -- (0.34,0);
    \foreach \x in {-0.26,-0.13,0,0.13,0.26} {
      \draw (\x,0) -- (\x+0.08,0.10);
    }
  \end{scope}
  \node[above left=4pt] at (1.20,3.70) {$\mathrm{M}_3$};

  % Telescope T
  \draw[fill=white] (1.16,0.36) -- (1.16,-0.12) arc (180:360:0.14 and 0.05) -- (1.44,0.36);
  \draw (1.30,0.36) ellipse (0.14 and 0.05);
  \draw (1.24,-0.16) -- (1.24,-0.32) arc (180:360:0.06 and 0.025) -- (1.36,-0.16);
  \node[left=3pt] at (1.16,0.02) {$\mathrm{T}$};
  \node[right=3pt, align=center] at (1.46,0.02) {望\\[-2pt]远\\[-2pt]镜};
\end{tikzpicture}

\end{center}
\end{problem}
\clearpage


\begin{problem}{12.1.3}
改用本题图的实验装置做上述菲佐实验, 半径为 R、折射率为 n 的玻璃圆柱绕固定的 O 点以恒定角速度 $\omega$ 转动. 两束光通过玻璃的距离均为 2l, 求射入望远镜 T 的

两束光的光程差,设圆柱拖动以太的系数为 k.
\begin{center}
\begin{tikzpicture}[>=Stealth, line width=0.7pt]
  % Glass cylinder cross-section of radius R = 1.65
  \draw (0,0) circle (1.65);
  \fill (0,0) circle (0.055);

  % Chord of length l at y = 0.85 (half-chord = sqrt(1.65^2 - 0.85^2) = 1.414)
  \draw (-1.414,0.85) -- (1.414,0.85);

  % Dimension indicator l above chord
  \draw[thin] (-1.414,1.05) -- (-1.414,1.45);
  \draw[thin] (1.414,1.05) -- (1.414,1.45);
  \draw[<->] (-1.414,1.26) -- (1.414,1.26) node[midway, fill=white, inner sep=2pt] {$l$};

  % Radii and perpendicular distance d
  \draw (0,0) -- (-1.414,0.85) node[midway, below left=1pt] {$R$};
  \draw (0,0) -- (0,0.85) node[midway, left=1pt] {$d$};
  \draw (0,0) -- (1.10,0.85) node[midway, below right=1pt] {$r$};
  \node[below right=1pt] at (1.08,0.85) {$A$};

  % Angle alpha
  \draw (0,0.42) arc (90:37.7:0.42);
  \node at (63.8:0.62) {$\alpha$};
\end{tikzpicture}

\end{center}
\end{problem}
\clearpage


\begin{problem}{12.1.4}
一个速度为 $v$ 的航天旅行者和他的地球上的朋友，在出发时对好了钟的时刻分别为 $t' = 0$ （旅行者的钟）和 $t = 0$ （地球上的钟）。地球上的朋友同时观察两个钟，直接观察 $t$ ，用望远镜观察 $t'$ ，试问，当 $t'$ 读数为1小时时， $t$ 的读数为多少？
\end{problem}
\clearpage


\begin{problem}{12.1.5}
在地球上某地发生的一个过程经历的时间为 $\Delta t$ ，问在以每小时 1800km 的恒定速度飞行的飞机上的观察者测得的此过程经历的时间为多大？若过程开始时两钟读数相同，经过多长时间飞机上的钟的读数又和地球上的钟的读数相同。
\end{problem}
\clearpage


\begin{problem}{12.1.6}
静止时 $\mu$ 子的平均寿命为 $2.197 \times 10^{-6} s$ ，在宇宙射线中， $\mu$ 子的速度可达 0.99c (c 为真空中的光速)，在地球上的观察者看来，平均来讲，这种高速的 $\mu$ 子能运动多大距离？
\end{problem}
\clearpage


\begin{problem}{12.1.7}
在离观测者距离为 $L$ 的位置上有一只时钟 $B$ , 说明怎样能使这只钟与观测者身边的时钟 $A$ 时间上一致?
\end{problem}
\clearpage


\begin{problem}{12.1.8}
设 $S'$ 系以速度 $\nu$ 相对于 $S$ 系运动, 在两个坐标原点重合的时刻, 放在两个坐标原点的时钟指示同一时刻, $t = t' = 0$ . 试问: 在 $S$ 系和 $S'$ 系中以后都有 $t = t'$ 的点 $(x, y, z)$ 和 $(x', y', z')$ 将分别随 $t$ 和 $t'$ 如何变化?
\end{problem}
\clearpage


\begin{problem}{12.1.9}
假如一艘宇宙飞船以恒定速度 $v = \sqrt{0.9999} c$ 飞行，试问飞船飞到距地球4光年远的半人马星座并返回地球的旅行中应储备够用多少时间的粮食及其他装备？
\end{problem}
\clearpage


\begin{problem}{12.1.10}
一列车以 $v = 0.6c$ ( $c$ 为真空中的光速)的恒定速度运动, 经过地面上 $A, B$ 两点所用的时间为 $40\mathrm{min}$ (以列车上静止的钟计时). 求:

(1) 地面上的人测得的 A、B 间的距离；

(2) 列车上的人测得的 A、B 间的距离；

(3) 在列车上的人看来,地面上 B 处的钟比 A 处的钟拨快了多少?
\end{problem}
\clearpage


\begin{problem}{12.1.11}
在惯性系 $S'$ 中的同一地点先后发生两个事件, 时间间隔为 $\Delta t' = 300\mathrm{s}$ , 在另一惯性系 $S$ 中测得这两个事件的时间间隔为 $\Delta t = 500\mathrm{s}$ . 求在惯性系 $S$ 中测得的这两个事件发生的地点相距多远?
\end{problem}
\clearpage


\begin{problem}{12.1.12}
静长为 L 的车厢以 v 的恒定速度运动, 在车厢的后端 A 沿运动方向发出一光信号, 经前端 B 的平面镜反射回到后端 A. 求:

(1) 在车厢里的人看来, 光信号从发出到到达 B 端所需的时间 $\Delta t_{1}^{\prime}$ 以及从 A 端发出到返回 A 端所需的时间 $\Delta t_{2}^{\prime}$ ;

(2) 在地面上的人看来, 与 $\Delta t_{1}^{\prime}, \Delta t_{2}^{\prime}$ 相应的时间.
\end{problem}
\clearpage


\begin{problem}{12.1.13}
两条静止长度均为 $L_{0}$ 的宇宙飞船 $A, B$ 沿平行的轨道以恒定的相对速率相向而行， $A$ 船上的观察者看到自己的船头先与对方的船尾相遇， $\Delta t$ 时间后，自己的船尾再与对方的船头相遇，试问：

(1) B 船上的观察者看到两船相遇的次序是怎样的？

(2) 两船的相对速度多大?

(3) 看到两船以大小相等、方向相反的速度运动的观察者看到两船两端相遇的次序如何？
\end{problem}
\clearpage


\begin{problem}{12.1.14}
一艘火箭飞船以 $v = 0.8c$ 的速度飞经地球, 相遇时都把时间调整到 12 点整.

（1）飞船上的时钟 12 点半时飞经一相对于地球静止、它的钟也与地球同步的宇航站，此时宇航站的指示何时？

(2) 地球上的人测得的宇航站离地球多远?

(3) 飞船经宇航站时, 向地球发一信号, 地球上的人何时收到这个信号?

(4) 地球上的人收到信号后立即发回信号, 飞船上的人何时接到信号?
\end{problem}
\clearpage


\begin{problem}{12.1.15}
一根米尺相对于 $S'$ 系静止，与 $o'x'$ 轴的夹角为 $30^\circ$ ，如 $S'$ 系相对于 $S$ 系沿 $x$ 轴（与 $x'$ 轴平行）方向运动，在 $S$ 系中测得米尺与 $x$ 轴的夹角为 $45^\circ$ ，求：

(1) S 系中测得的米尺的长度；

(2) $S'$ 系相对于S系的速度.
\end{problem}
\clearpage


\begin{problem}{12.1.16}
一短跑选手在地球上以 $10\mathrm{s}$ 的时间跑完 $100\mathrm{m}$ , 在沿跑的方向以 $0.98c$ 的速度相对于地球飞行的飞船中的观察者看到的选手跑了多少时间和多长距离?
\end{problem}
\clearpage


\begin{problem}{12.1.17}
若 12.1.11 题 $S'$ 系中先后发生的两个事件并不发生在同一地点, 在有相对运动的方向上相距 $\Delta x' = 10$ 光分, $\Delta t'$ 、 $\Delta t$ 仍和 12.1.11 题相同. 求:

(1) $S'$ 系相对于 $S$ 系的速度 $v$ ;

(2) 在 S 系中测得的两事件发生的地点间的距离 $\Delta x$ .
\end{problem}
\clearpage


\begin{problem}{12.1.18}
一根杆自左向右运动, 当这杆的左端经过一架照相机时, 这根杆和一根静止的标有刻度的米尺一起拍摄在一张照片上, 照片显示: 杆的左端与米尺上的 0 标记重合, 右端与米尺的 $0.9 \mathrm{~m}$ 标记重合. 如果杆是以 $v = 0.8 c$ 的速度相对于照相机运动的, 求杆的实际长度.
\end{problem}
\clearpage


\begin{problem}{12.1.19}
像上题那样, 在一根 1m 长的杆的中点经过照相机的瞬间, 打开照相机快门, 连同一根静止的有刻度的米尺拍下照片, 如果杆相对于照相机的速度 v=0.8c, 照片上记录下来的运动杆的长度该是多少?
\end{problem}
\clearpage


\begin{problem}{12.1.20}
一个粒子在 $x^{\prime}y^{\prime}$ 平面内以 $v = \frac{c}{2}$ 的恒定速度在 $t' = 0$ 时离开 $O'$ 点， $\bm{v}$ 与 $x'$ 轴的夹角为 $60^{\circ}, O'$ 点沿 $x'$ 亦沿 $x$ 轴相对于 $O$ 的速度为 0.6c. $t = 0$ 时， $O'$ 、 $O$ 点重合，试求由 $O$ 所确定的粒子的运动方程.
\end{problem}
\clearpage


\begin{problem}{12.1.21}
观测者 $O$ 和 $O'$ 以 $0.6c$ 的相对速度互相接近, 如果 $O$ 测得 $O'$ 离他的初始距离为 $20m$ , 按 $O$ 的测定, 多长时间后与 $O'$ 相遇? 按 $O'$ 的测定, 多长时间后与 $O$ 相遇?
\end{problem}
\clearpage


\begin{problem}{12.1.22}
一个还能活60年的人打算访问距离我们160000光年的遥远星系，他的恒定速度必须多大？
\end{problem}
\clearpage


\begin{problem}{12.1.23}
一枚火箭飞船相对于地球的速度为 $0.8c(c$ 为真空中的光速), $O'$ 和 $O$ 分别是飞船和地球上的观测者, $x = x' = 0$ 时 $t = t' = 0$ . $O$ 用望远镜看 $O'$ 的时候, 他自己的钟读数为 $30s$ , 问他看到 $O'$ 的时钟读数是多少?
\end{problem}
\clearpage


\begin{problem}{12.1.26}
A、B 两钟相对于某惯性系的一条直线上做方向相同、速率不同的匀速运动。在相遇时两钟均拨到零点，A 钟在时刻 $T_{A}$ 向 B 的方向发出一光信号，B 钟收到该信号的时间为 $T_{B}$ 。证明：B 钟相对于 A 钟的运动速率为

\[
v = \frac {c (T _ {B} ^ {2} - T _ {A} ^ {2})}{(T _ {A} ^ {2} + T _ {B} ^ {2})}
\]
\end{problem}
\clearpage


\begin{problem}{12.1.27}
一质点在惯性系 $S'$ 中做匀速率圆周运动，轨迹方程为

\[
x ^ {\prime 2} + y ^ {\prime 2} = r ^ {2}, \quad z ^ {\prime} = 0
\]

$S'$ 系以速度 $v$ 沿 $x$ 轴向正向运动，在 $t = t' = 0$ 时， $x, x', y, y', z, z'$ 轴分别重合。试证明：在惯性系 $S$ 中测得的轨道为一椭圆，椭圆的中心以速度 $v$ 移动。
\end{problem}
\clearpage


\begin{problem}{12.1.28}
在 O 的参考系中有一个静止的正方形, 其面积为 $100 ~m^{2}$ , 观察者 $O'$ 以 0.8c 的速度沿正方形的对角线运动. $O'$ 测得的面积多大?
\end{problem}
\clearpage


\begin{problem}{12.1.29}
在 12.1.27 题中所述的对 $S'$ 系静止的圆, 在 $S$ 系中的观测者测得其图形面积多大?
\end{problem}
\clearpage


\begin{problem}{12.1.30}
静止长度为 $l_{0}$ 的两根尺, 平行于 x 轴相向做匀速运动, 与一根尺相连结的观察者发现,两根尺的左端和右端两次重合的时间间隔为 $\Delta t$ ,试问两根尺的相对速度 v 多大?
\begin{center}
\begin{tikzpicture}[>=Stealth, line width=0.7pt]
  % Moving contracted ruler (top)
  \draw (0,0.82) rectangle (3.22,1.04);
  \draw[->] (3.32,0.93) -- (4.15,0.93) node[midway, above=2pt] {$v$};

  % Top length dimension (clear of top ruler edge)
  \draw[thin] (0,1.22) -- (0,1.92);
  \draw[thin] (3.22,1.22) -- (3.22,1.92);
  \draw[<->] (0,1.58) -- (3.22,1.58)
    node[midway, fill=white, inner sep=2pt] {$l_0\sqrt{1-\left(\frac{v}{c}\right)^2}$};

  % Rest ruler (bottom)
  \draw (0,0) rectangle (4.60,0.22);

  % Bottom length dimension
  \draw[thin] (0,-0.15) -- (0,-0.55);
  \draw[thin] (4.60,-0.15) -- (4.60,-0.55);
  \draw[<->] (0,-0.35) -- (4.60,-0.35) node[midway, fill=white, inner sep=3pt] {$l_0$};
\end{tikzpicture}

\end{center}
\end{problem}
\clearpage


\begin{problem}{12.1.31}
静止半径为 $R_0$ 的球, 相对于远处的观察者以速度 $\bm{v}$ 运动, 球上有明显的标记, 当观察者看到球速正好与他和球的连线垂直时, 拍了张照片, 当他冲洗胶片时看到了什么?
\end{problem}
\clearpage


\begin{problem}{12.1.32}
一个原子钟被喷气式飞机带着绕地球一

周后,与没有动的精确同步的同样的钟比较,适当近似,狭义相对论预言有多大的不一致?
\end{problem}
\clearpage


\begin{problem}{12.1.33}
设 $S'$ 系以速度 $v$ 沿 $x$ 轴正方向相对于 $S$ 系运动, 固定在 $S'$ 系内 $(x_0', y_0', z_0')$ 点的时钟, 在 $t'_0$ 时刻从 $S$ 系内 $(x_0, y_0, z_0)$ 点旁边经过, 这时 $S$ 系内的时钟指示 $t_0$ 时刻, 写出在这种情况下的洛伦兹变换公式.
\end{problem}
\clearpage


\begin{problem}{12.1.34}
设 $S'$ 系相对于 $S$ 系的速度 $V$ 并不平行于 $x$ 轴, $t = t' = 0$ 时, $x$ 与 $x'$ 轴, $y$ 与 $y'$ 轴、 $z$ 与 $z'$ 轴两两重合. 试导出这种情况下的洛伦兹变换公式.

提示: 把 $r, r'$ 分解为平行于 V 的分量 $r_{\parallel} = \frac{r \cdot V}{v^{2}} V, r'_{\parallel} = \frac{r' \cdot V}{v^{2}} V$ 和垂直于 V 的分量 $r_{\perp} = r - r_{\parallel}, r'_{\perp} = r' - r'_{\parallel}$ .
\end{problem}
\clearpage


\begin{problem}{12.1.35}
考虑图 12.7 所示的三个坐标系 $S, S'$ 和 $S''$ 之间的洛伦兹变换, 这里 $x, x'$ , $x''$ 是互相平行的, $S', S''$ 沿 $x$ 轴向正向运动, $S'$ 相对于 $S$ 的速度为 $v_1, S''$ 相对于 $S'$ 的速度为 $v_2$ . 证明: 洛伦兹变换的逆变换是洛伦兹变换, 两个洛伦兹变换的结果是另一个洛伦兹变换.
\begin{center}
\begin{tikzpicture}[>=Stealth, line width=0.7pt]
  % Frame S
  \draw[->] (0,0) -- (1.35,0) node[below=2pt] {$x$};
  \draw[->] (0,0) -- (0,1.35) node[right=2pt] {$y$};
  \draw[->] (0,0) -- (-0.48,-0.48) node[below right=1pt] {$z$};
  \node at (0.95,1.08) {$S$};

  % Frame S'
  \begin{scope}[shift={(2.15,0)}]
    \draw[->] (0,0) -- (1.35,0) node[below=2pt] {$x'$};
    \draw[->] (0,0) -- (0,1.35) node[right=2pt] {$y'$};
    \draw[->] (0,0) -- (-0.48,-0.48) node[below right=1pt] {$z'$};
    \node at (0.95,1.08) {$S'$};
    \draw[->] (0,0.72) -- (0.85,0.72);
    \node[below right=2pt and 3pt] at (0,0.72) {$v_1(\mbox{对}\,S)$};
  \end{scope}

  % Frame S''
  \begin{scope}[shift={(4.30,0)}]
    \draw[->] (0,0) -- (1.35,0) node[below=2pt] {$x''$};
    \draw[->] (0,0) -- (0,1.35) node[right=2pt] {$y''$};
    \draw[->] (0,0) -- (-0.48,-0.48) node[below right=1pt] {$z''$};
    \node at (0.95,1.08) {$S''$};
    \draw[->] (0,0.72) -- (0.85,0.72);
    \node[below right=2pt and 3pt] at (0,0.72) {$v_2(\mbox{对}\,S')$};
  \end{scope}
\end{tikzpicture}

\end{center}
\end{problem}
\clearpage


\section{狭义相对论的运动学}

\begin{problem}{12.2.1}
$S'$ 系以 $V$ 的恒定速度沿 $S$ 系的 $x$ 轴向正向运动, 若质点在 $S'$ 系的 $x'y'$ 平面内运动, 速度 $v'$ 与 $x'$ 轴的夹角为 $\theta'$ . 试证明, 在 $S$ 系中, 质点的速度 $v$ 与 $x$ 轴的夹角 $\theta$ 满足关系:

\[
\tan \theta = \frac {v ^ {\prime} \sin \theta^ {\prime} \sqrt {1 - \left(\frac {V}{c}\right) ^ {2}}}{V + v ^ {\prime} \cos \theta^ {\prime}}
\]
\end{problem}
\clearpage


\begin{problem}{12.2.2}
两根静止长度均为 $l_{0}$ 的杆成一直线放置，并以相对于一个参考系以相同的速率 $\bm{v}$ 匀速相向运动。试问在固连于一个杆的参考系内的观察者测得的另一根杆的长度是多少？
\end{problem}
\clearpage


\begin{problem}{12.2.3}
两根静止长度均为 $l_{0}$ 的尺在某一参考系中分别如图12.8(a)、(b)那样以相同的速率 $v$ 、互成 $90^{\circ}$ 角运动，在与一根尺固连的参考系中测得的另一根尺的长度是多少？
\begin{center}
\begin{tikzpicture}[>=Stealth, line width=0.8pt]
    % Vertical ruler moving upward
    \draw (0, 1.0) rectangle (0.24, 3.4);
    \draw[->] (0.58, 1.7) -- (0.58, 2.7) node[midway, right=2pt] {$v$};

    % Horizontal ruler moving rightward
    \draw (-0.45, 0.25) rectangle (1.95, 0.49);
    \draw[->] (2.10, 0.37) -- (3.10, 0.37) node[midway, above=2pt] {$v$};

    % Subfigure label
    \node at (0.75, -0.15) {(b)};
\end{tikzpicture}

\end{center}
\end{problem}
\clearpage


\begin{problem}{12.2.4}
在某一参考系中, 两条飞船 A、B 分别以 0.8c 及 0.6c 的速率相向而行, 找出一个参考系 $S'$ , 使两条飞船分别以相同的速率 $v'$ 相向而行. 求:

(1) $S'$ 系对 $S$ 系的速度 $\bm{v}$ 

(2) $v'$ 有多大?
\end{problem}
\clearpage


\begin{problem}{12.2.5}
在某一参考系中, 物体 $A$ 以匀速率 $v_{A} = 0.8c$ 沿 $x$ 轴向正向运动, 物体 $B$ 以匀速率 $v_{B} = 0.6c$ 沿 $x$ 轴向负向运动. 在 $A$ 看来, $B$ 的速度 $v_{B}'$ 如何?
\end{problem}
\clearpage


\begin{problem}{12.2.6}
利用速度变换公式, 求出 12.1.2 题所述的菲佐实验中, 顺水速和逆水速两种情况下的光速, 与用以太假设解释菲佐实验引入的拖曳系数 k 后的光速相比较, 从而导出 k 与折射率的关系.
\end{problem}
\clearpage


\begin{problem}{12.2.7}
一束光通过折射率为 n 的玻璃块, 如果玻璃块以恒定速度 v 沿入射光方向运动. 试求在玻璃块中光相对于实验室参考系的速度.
\end{problem}
\clearpage


\begin{problem}{12.2.8}
火箭 A 相对于地球以 0.8c 向 +y 方向飞行, 火箭 B 以 0.6c 向 -x 方向飞行, 由火箭 B 测得的火箭 A 的速度是多少?
\end{problem}
\clearpage


\begin{problem}{12.2.9}
一个处于静止状态的 $\mathbf{K}^{\circ}$ 介子衰变成一个 $\pi^{+}$ 介子和一个 $\pi^{-}$ 介子, 各自都具有 $0.827c$ 的速率, 如果 $\mathbf{K}^{\circ}$ 介子在以 $0.6c$ 的速度飞行时衰变, 其中一个 $\pi$ 介子可能有的最大速率多大?
\end{problem}
\clearpage


\begin{problem}{12.2.10}
证明 $v = \frac{\sqrt{(v' + V)^2 - (v' \times V)^2 / c^2}}{1 + v' \cdot V / c^2}$ ，其中 $v, v'$ 分别是粒子在 $S$ 系和 $S'$ 系中的速度， $V$ 是 $S'$ 系相对于 $S$ 系的速度。
\end{problem}
\clearpage


\begin{problem}{12.2.11}
当 $S'$ 系相对于 S 系的速度 v 有任意方向时, 试导出速度变换公式, 并表示成矢量形式.
\end{problem}
\clearpage


\begin{problem}{12.2.12}
一光源在其静止的坐标系中向各方向均匀地辐射, 当它以接近于光速 $c$ 的速率 $v$ 在参考系 $S$ 中运动时, 在 $S$ 系中, 它的辐射将在它的前进方向上有强烈的辐射 (这种现象称为前灯效应). 要使发射出的辐射有一半在半顶角 $\theta = 10^{-3} \mathrm{rad}$ 的圆锥状立体角内, $v$ 应多大?

提示: 可利用 12.2.1 题结果.
\end{problem}
\clearpage


\begin{problem}{12.2.13}
证明光线与 $x$ 轴的夹角 $\theta$ 的变换公式为

\[
\tan \theta = \frac {c \sin \theta^ {\prime} \sqrt {1 - \left(\frac {V}{c}\right) ^ {2}}}{c \cos \theta^ {\prime} + V}
\]

或

\[
\cos \theta = \frac {c \cos \theta^ {\prime} + V}{c + V \cos \theta^ {\prime}}
\]
\end{problem}
\clearpage


\begin{problem}{12.2.14}
一光束在某一参考系内构成一立体角 $\mathrm{d}\Omega$ ，试问当变换到另一个惯性参考系时，其立体角如何变化？
\end{problem}
\clearpage


\begin{problem}{12.2.15}
相对于有限的恒星系静止的观察者 A, 看到恒星系有各向同性的星体分布, 若 A 能看到的全部星体数为 N, 则在立体角 $d\Omega$ 中看到的星体数为

\[
\mathrm{d} N (\theta , \varphi) = \frac {N}{4 \pi} \mathrm{d} \Omega (\theta , \varphi)
\]

其中 $\mathrm{d}\Omega (\theta ,\varphi) = \sin \theta \mathrm{d}\theta \mathrm{d}\varphi$ 是 $\theta \sim \theta +\mathrm{d}\theta ,\varphi \sim \varphi +\mathrm{d}\varphi$ 所张的立体角.

另一个在沿 $z$ 轴以速度 $\bm{v}$ 运动的参考系中的观察者 $B$ 如能看到观察者 $A$ 看到的每一个星体. 求：

(1) 观察者 B 在 $\theta'$ 、 $\varphi'$ 的单位立体角内看到的星体数 $N(\theta', \varphi')$ ;

(2) 当 $v \to c$ 时, 观察者 $B$ 在 $\theta' = 0$ 和 $\theta' = \pi$ 处的单位立体角内看到的星体数 $N(0, \varphi')$ 、 $N(\pi, \varphi')$ .
\end{problem}
\clearpage


\begin{problem}{12.2.16}
宇宙飞船以 v=0.8c 的速率飞离地球，当在地球参考系中测得距地球 $6.66 \times 10^{8} km$ 时，一无线电信号由地球上的观察者向飞船发出.

(1) 在飞船参考系中测量多长时间到达飞船?

(2) 在地球参考系中测量多长时间到达飞船?

(3) 在地球参考系中测量飞船收到信号时的位置.
\end{problem}
\clearpage


\begin{problem}{12.2.17}
在一惯性系 $S$ 中观察到两宇宙飞船沿两直线相向平行飞行, 轨道间距为 $d$ , 如图 12.9 所示. 每条飞船的速率均为 $\frac{c}{2}, c$ 为真空中的光速.

(1) 在 S 系中看, 当两飞船抵达最近点(图中虚线表示的 A、B 两点)时, 飞船(1)以 $\frac{3}{4}c$ 的速率投出一小包裹. 从飞船(1)上的观察者来看, 为使包裹能被飞船(2)收到, 必须以什么样的角度投出?

(2) 从飞船(1)的观察者来看,包裹的速率多大?
\begin{center}
\begin{tikzpicture}[>=Stealth, line width=0.8pt]
    % Parallel trajectories
    \draw[dashed] (0, -1.65) -- (0, 1.65);
    \draw[dashed] (2.6, -1.65) -- (2.6, 1.65);

    % Spaceship (1) moving downward
    \draw[fill=white] (-0.11, 1.48) -- (0.11, 1.48)
        .. controls (0.11, 1.18) and (0.04, 1.02) .. (0, 0.96)
        .. controls (-0.04, 1.02) and (-0.11, 1.18) .. (-0.11, 1.48) -- cycle;
    \node[right=4pt] at (0.11, 1.30) {$(1)$};
    \draw[->] (-0.38, 1.52) -- (-0.38, 0.78) node[midway, left=2pt] {$\frac{c}{2}$};

    % Spaceship (2) moving upward
    \draw[fill=white] (2.49, -1.48) -- (2.71, -1.48)
        .. controls (2.71, -1.18) and (2.64, -1.02) .. (2.6, -0.96)
        .. controls (2.56, -1.02) and (2.49, -1.18) .. (2.49, -1.48) -- cycle;
    \node[left=4pt] at (2.49, -1.30) {$(2)$};
    \draw[->] (2.98, -1.52) -- (2.98, -0.78) node[midway, right=2pt] {$\frac{c}{2}$};

    % Closest-approach dashed line AB with end ticks
    \draw[dashed] (0, 0) -- (2.6, 0);
    \draw (0, -0.07) -- (0, 0.07);
    \draw (2.6, -0.07) -- (2.6, 0.07);
    \node[left=3pt] at (0, 0) {$A$};
    \node[right=3pt] at (2.6, 0) {$B$};

    % Track separation d
    \node (dlabel) at (1.3, -0.72) {$d$};
    \draw[->] (dlabel.west) -- (0.04, -0.72);
    \draw[->] (dlabel.east) -- (2.56, -0.72);
\end{tikzpicture}

\end{center}
\end{problem}
\clearpage


\begin{problem}{12.2.18}
观察者 1 看到粒子以速度 $v$ 在与他的 $z$ 轴成 $\varphi$ 角的直线轨道上运动, 观察者 2 以速度 $V$ 相对于观察者 1 沿 $z$ 方向运动, 求观察者 2 测得的粒子运动的速度, 并考虑 $v \to c$ 的极限情况, 检验你得到的结果.
\end{problem}
\clearpage


\begin{problem}{12.2.19}
一枚火箭相对于固连于自己的参考系以匀加速度 $a_0$ 做直线运动, 按照地球的钟, 从静止出发, 经过多长时间 (假定时钟的运行与其加速度无关), 火箭加速到 $V = 0.8c$ ? 按照火箭内的时钟, 加速到这一速度需用多长时间?
\end{problem}
\clearpage


\begin{problem}{12.2.20}
静止的氢原子发出一条光谱线 $H_{\delta}$ 的波长 $\lambda=0.4101\times10^{-6}m$ . 当氢原子以 $v=5\times10^{5}m/s$ 向着观察者飞行时, 观察者看到的波长有多大的改变?
\end{problem}
\clearpage


\begin{problem}{12.2.21}
当一光源以速率 $V_{1}$ 向地球靠近时, 地球上静止的人甲看到光源发出的是波长 $\lambda_{1}=5000\mathring{A}$ 的绿光, 而相对地球以速率 $V_{2}$ 运动 (与光源的运动沿同一直线) 的人乙看到的光是波长 $\lambda_{2}=6000\mathring{A}$ 的红光. 当光源以相同的速率 $V_{1}$ 远离地球时, 甲也看到波长 $\lambda_{2}$ 的红光.

(1) 求光谱发出的光的波长 $\lambda_{0}$ ,

(2) 求 $V_{1}$ 、 $V_{2}$ 的值；

(3) 当光源以 $V_{1}$ 远离地球时，乙看到的光的波长多大？
\end{problem}
\clearpage


\begin{problem}{12.2.22}
一辆汽车以每小时 $120\mathrm{km}$ 的速度接近一个汽车速度监视雷达站，如果雷达站的工作频率 $\nu_0 = 2\times 10^{10}\mathrm{Hz}$ ，雷达站上的警察测到的频移有多大？
\end{problem}
\clearpage


\begin{problem}{12.2.23}
假定眼睛可以看到的最大波长是 $6500 \, \AA$ ，要使火箭发出的绿光 ( $5000 \, \AA$ )，对于地球上的观测者成为不可见的，该火箭必须以多大速度运动？
\end{problem}
\clearpage


\begin{problem}{12.2.24}
氢原子光谱中的一条谱线 $\mathbf{H}_{\alpha}$ , 波长为 $6.561 \times 10^{-7} \mathrm{~m}$ , 在地球上测得来自太阳赤道上相对两端辐射 $\mathbf{H}_{\alpha}$ 线的波长差为 $9 \times 10^{-12} \mathrm{~m}$ . 假定此效应是由于太阳自转引起的, 求太阳的自转周期 $T$ , 已知太阳直径 $D = 1.4 \times 10^{9} \mathrm{~m}$ .
\end{problem}
\clearpage


\begin{problem}{12.2.25}
每个双生子的心脏一秒钟跳一次,每人每次心跳时发出一个无线电脉冲,留在家里的一位在其惯性系中保持静止,另一位旅行者在零时刻由静止出发,极快地加速到速度 $v$ (在比一次心跳更短的时间内,且不扰乱他的心脏),旅行者用他的钟记录自己旅行了时间 $t_1$ ,同时一直发出脉冲,也一直接收从家里发来的脉冲.在时刻 $t_1$ ,他突然改变速度方向,并于时刻 $2t_1$ 回到家中,他共计发出多少脉冲?在去途中他收到多少脉冲?归途中收到多少脉冲?总计收到和发出的脉冲之比为何值?家里的那位从时刻零到 $t_2$ (用他的钟计时)收到多普勒降频脉冲,从时刻 $t_2$ 开始收到多普勒升频脉冲,设 $t_3$ 是从时刻 $t_2$ 直至旅行结束的时间间隔,他在时间间隔 $t_2$ 内收到多少脉冲?在 $t_3$ 内又收到多少?两者之比为何值?发出和接收的脉冲总数之比又为何值?
\end{problem}
\clearpage


\begin{problem}{12.2.26}
一宇宙飞船以速度 $0.6c$ 离开地球飞向离地球3光年的宇航站（地球参考系测出的距离），然后再返回地球。起飞时地球的钟为 $t = 0$ ，飞船的钟为 $t' = 0$ 。地球上每隔1年向飞船发一光脉冲，飞船上的人也每隔1年向地球发一光脉冲，试问：

(1) 旅行过程中,地球上的人和飞船上的人各发了几个脉冲?

(2) 飞船上的人在哪些时刻收到脉冲？

(3) 地球上的人在哪些时刻收到脉冲？

(1) 按在地球参考系(S 系)中测量,宇航站离地球的距离为 L=3c (时间以年为单位).

飞船从地球到宇航站需时 $t_1 = \frac{L}{0.6c} = 5$ （年），从宇航站返回地球需时 $t_2 = \frac{L}{0.6c} = 5$ （年). 整个旅行过程，需时 $t_1 + t_2 = 10$ 年. 地球上的人向飞船共发了9个脉冲（回到地球就

不再互发脉冲).

对飞船到达宇航站这个事件用洛伦兹变换公式,求飞船的钟的示数,

\[
t = \frac {t ^ {\prime} + \frac {v}{c ^ {2}} x ^ {\prime}}{\sqrt {1 - \left(\frac {v}{c}\right) ^ {2}}}
\]

今 $v = 0.6c, x' = 0, t = t_1 = 5$ 年.

\[
t ^ {\prime} = t \sqrt {1 - \left(\frac {v}{c}\right) ^ {2}} = 4 \text {年}
\]

按飞船的钟计,从宇航站返回地球又需4年,因此旅行期间,飞船上的人向地球共发了7个脉冲.

(2) 方法一: 按地球上的人测量, 地球上一年后发第一个脉冲, 飞船上的人收到时刻为 $t_1$ (以年计).

\[
\begin{array}{c} {c (t _ {1} - 1) = 0. 6 c t _ {1}} \\ {t _ {1} = 2. 5 (\text {年})} \end{array}
\]

用洛伦兹变换公式,可算出,飞船上的人收到第一个脉冲的时刻为

\[
t _ {1} ^ {\prime} = t _ {1} \sqrt {1 - \left(\frac {v}{c}\right) ^ {2}} = 2 (\text {年})
\]

飞船上的人收到来自地球上的人发来的第二个脉冲的时刻 $t_2$ (用 $S$ 系) 和 $t_2'$ (用 $S'$ 系) 计算如下：

\[
\begin{array}{r l} & {c (t _ {2} - 2) = 0. 6 c t _ {2}, \quad t _ {2} = 5 \text {年}} \\ & {\qquad t _ {2} ^ {\prime} = t _ {2} \sqrt {1 - \left(\frac {v}{c}\right) ^ {2}} = 4 \text {年}} \end{array}
\]

飞船收到第三个光脉冲是在归程中收到的. 飞船收到时已走过的距离加上第三个脉冲走过的距离等于往返的整个路程.(按 $S$ 系测出)

\[
0. 6 c \cdot t _ {3} + c (t _ {3} - 3) = 2 \times 3 c
\]

飞船上的人收到第三个脉冲时刻( $S'$ 系测得)

\[
t _ {3} ^ {\prime} = t _ {3} \sqrt {1 - \left(\frac {v}{c}\right) ^ {2}} = \frac {9}{1 . 6} \times 0. 8 = 4. 5 (\text {年})
\]

飞船上的人收到第 $n$ 个 $(n\geqslant 3)$ 脉冲，

\[
\begin{array}{r l} & c (t _ {n} - n) + 0. 6 c t _ {n} = 2 \times 3 c \\ & \quad 1. 6 c t _ {n} = (n + 6) c \end{array}
\]

\[
t _ {n} ^ {\prime} = t _ {n} \sqrt {1 - \left(\frac {v}{c}\right) ^ {2}} = \frac {n + 6}{1 . 6} \times 0. 8 = \frac {1}{2} (n + 6) (\text {年}) \qquad (3 \leqslant n \leqslant 9)
\]

飞船上的人收到地球发来的9个脉冲的时刻(飞船的钟)为 $t'=2,4,4.5,5,5.5,6,6.5,7$ 和7.5年.

(3) 方法一: 飞船上发第 1 个脉冲时 $t_{10}^{\prime}=1$ 年, 此时地球上的钟

\[
t _ {1 0} = \frac {t _ {1 0} ^ {\prime} + \frac {v}{c ^ {2}} x _ {1 0} ^ {\prime}}{\sqrt {1 - \left(\frac {v}{c}\right) ^ {2}}} = \frac {1}{0 . 8} = 1. 2 5 (\text {年})
\]

地球上收到第1个脉冲，地球上的钟示数为 $t_{1}$ ，

\[
c (t _ {1} - t _ {1 0}) = v t _ {1 0}
\]

\[
t _ {1} = t _ {1 0} \Big (1 + \frac {v}{c} \Big) = 2 (\text {年})
\]

同样考虑在飞船返回以前发的第 $n$ 个脉冲被地球收到的时刻 $t_n$ 

\[
c (t _ {n} - t _ {n 0}) = v t _ {n 0}
\]

\[
t _ {n} = \left(1 + \frac {v}{c}\right) t _ {n 0} = \left(1 + \frac {v}{c}\right) \frac {t _ {n 0} ^ {\prime}}{\sqrt {1 - \left(\frac {v}{c}\right) ^ {2}}} = \sqrt {\frac {1 + \frac {v}{c}}{1 - \frac {v}{c}}} \cdot n = 2 n
\]

其中用了 $t_{n0}^{\prime} = n, n = 1,2,3,4.$ 

飞船上发第5个脉冲时， $t_{50}^{\prime} = 5$ 年，此时地球上的钟为

\[
t _ {5 0} = \frac {t _ {5 0} ^ {\prime}}{\sqrt {1 - \left(\frac {v}{c}\right) ^ {2}}}
\]

地球上的人收到第5个脉冲的时刻 $t_5$ 

\[
c \left(t _ {5} - t _ {5 0}\right) + v t _ {5 0} = 6 c
\]

\[
t _ {5} = t _ {5 0} \left(1 - \frac {v}{c}\right) + 6 = t _ {5 0} ^ {\prime} \sqrt {\frac {1 - \frac {v}{c}}{1 + \frac {v}{c}}} + 6 = 8. 5 (\text {年})
\]

同样可得

\[
t _ {n} = \frac {1}{2} n + 6 (\text {年}) \qquad n = 5, 6, 7
\]

地球上的人收到7个脉冲的时间分别为

\[
t = 2, 4, 6, 8, 8. 5, 9, 9. 5 (\text { 年 })
\]

方法二:用多普勒效应来回答(2)、(3)问.

设地球发射脉冲的频率为 $\nu$ ，则当飞船远离地球时，飞船收到脉冲的频率 $\nu$ 为

\[
\nu^ {\prime} = \sqrt {\frac {c - v}{c + v}} \nu
\]

发射的周期和接收的周期分别用 T 和 $T'$ 表示, T=1 年, 则

\[
T ^ {\prime} = \frac {1}{\nu^ {\prime}} = \sqrt {\frac {c + v}{c - v}} \frac {1}{\nu} = \sqrt {\frac {c + v}{c - v}} T = 2 \text {年}
\]

飞船收到第1个脉冲时, $t_{1}^{\prime}=2$ 年,收到第2个脉冲时, $t_{2}^{\prime}=4$ 年.此时,由(1)问的解答知,飞船开始返航.

当飞船返航接近地球时，

\[
T ^ {\prime} = \sqrt {\frac {c - v}{c + v}} T = 0. 5 T = 0. 5 \text {年}
\]

因此，飞船收到第3、4、…、9个脉冲的时刻分别为 $t_{3}^{\prime}=4.5$ 年， $t_{4}^{\prime}=5$ 年， $\cdots,t_{9}^{\prime}=7.5$ 年.

同样考虑飞船发的脉冲周期为1年，当飞船远离地球时，地球上接收到的脉冲周期为2年，远离时飞船共发了4个脉冲，地球上收到这4个脉冲的时间分别为2年、4年、6年、8年。飞船返回期间发的脉冲，地球上接收到的周期为0.5年，因此地球上收到第5、6、7三个脉冲的时间分别为8.5年、9年、9.5年（均以地球的钟计时）。
\end{problem}
\clearpage


\begin{problem}{12.2.27}
证明光的多普勒效应公式：

（1）当光源和观察者相趋近时，观察者测得的频率 $\nu = \nu_0\sqrt{\frac{1 + \beta}{1 - \beta}}$ ，其中 $\beta = \frac{v}{c}, v$ 是相趋近时的相对速率. $\nu_0$ 是光的固有频率，即在光源为静止的参考系中测得的光的频率；

(2) 当光源和观察者相远离时, $\nu=\nu_{0}\sqrt{\frac{1-\beta}{1+\beta}}$ ;

(3) 当光源和观察者间的相对运动不在两者的连线上时, 观察者测得的频率为

\[
\nu = \sqrt {1 - \left(\frac {v}{c}\right) ^ {2}} \frac {c}{c - v \cos \theta} \nu_ {0}
\]

其中 $\theta$ 是在光源对观察者的速度 v 与光源、观察者连线的夹角.
\end{problem}
\clearpage


\begin{problem}{12.2.28}
一个光源静止在参考系 S 的原点 x=0，发射出两个光脉冲 $P_{1}$ 和 $P_{2}$ ，分别于 t=0 和 $t=\tau$ 发射，参考系 $S'$ 以 vi 相对于 S 系运动， $S'$ 系中的观察者在 $x'=0, t'=0$ 收到 $P_{1}$ .

(1) 求在 $x' = 0$ 点收到 $P_{2}$ 的时间 $\tau'$ ;

(2) 若 $\lambda, \lambda'$ 分别是在 $S$ 和 $S'$ 系中测量的光在真空中的波长，求 $\lambda'$ 与 $\lambda, v$ 的关系；

(3) 计算 $\mathrm{H}_{\beta}$ 辐射 $(\lambda = 4861.33\AA)$ 的 $\frac{v}{c}$ 的一级和二级多普勒频移. 这频移是质子通过 $20\mathrm{kV}$ 的电势差加速后辐射与平衡质子的辐射的频率差、假设加速后质子以保持不变的速度离去, 光谱计的光轴与质子的运动方向共线.
\end{problem}
\clearpage


\begin{problem}{12.2.29}
如图 12.13, 一频率为 $\nu$ 的单色横波在波源参考系 $S$ 中沿与 $x$ 方向成 $60^{\circ}$ 角传播, 该波源以速度 $v = 0$ . 8c 在 $x$ 方向向着在 $S'$ 系 (它的 $x'$ 轴平行于 $x$ 轴) 中静止的观察者运动. 求:

(1) 观察者测得的波的频率 $\nu$ ;

(2) 在 $S'$ 系中观测光的传播方向与 $x'$ 轴的夹角 $\theta'$ .
\begin{center}
\begin{tikzpicture}[>=Stealth, line width=0.8pt]
    % Frame S axes
    \draw[->] (0, 0) -- (3.8, 0) node[below=2pt] {$x$};
    \draw[->] (0, 0) -- (0, 2.2) node[left=2pt] {$y$};
    \node at (0.55, 1.45) {$S$};

    % Light source and propagation ray at 60 deg
    \coordinate (src) at (1.05, 0);
    \coordinate (Op) at (2.35, 2.25);
    \node[below=3pt] at (src) {\small 光源};
    \draw (src) -- (Op);

    % 60 degree angle arc
    \draw (src) ++(0:0.55) arc (0:60:0.55);
    \node at ($(src) + (28:0.86)$) {$60^\circ$};

    % Frame S' axes at O'
    \draw[->] (Op) -- ++(1.9, 0) node[below=2pt] {$x'$};
    \draw[->] (Op) -- ++(0, 1.75) node[left=2pt] {$y'$};
    \node at ($(Op) + (0.65, 1.35)$) {$S'$};
    \node[below right=1pt] at (Op) {$O'$};
    \node at ($(Op) + (0.55, -0.62)$) {\small 观察者};

    % Velocity v of observer/frame S'
    \draw[->] ($(Op) + (-0.15, 0.12)$) -- ++(-0.95, 0) node[midway, above=2pt] {$v$};
\end{tikzpicture}

\end{center}
\end{problem}
\clearpage


\begin{problem}{12.2.30}
一单色点光源发出频率为 $\nu$ 的辐射, 一观察者以匀速率 $v$ 沿一离光源最近距离为 $d$ 的直线运动, 如图 12.14 所示:

(1) 导出被观察到的频率的表达式, 将频率 $\nu$ 表示成图中 x 的函数;

(2) 若 $\frac{v}{c}=0.8$ ，画出(1)问的解答的近似图.
\begin{center}
\begin{tikzpicture}[>=Stealth, line width=0.8pt]
    % Axes for S and S'
    \draw[->] (0, 0) -- (4.8, 0) node[below=2pt] {$x, x'$};
    \draw[->] (0, 0) -- (0, 2.55);
    \node[below=2pt] at (0, 0) {$O$};

    % Light source on y-axis
    \coordinate (src) at (0, 2.0);
    \fill (src) circle (2pt);
    \node[above right=1pt and 2pt] at (src) {\small 光源};

    % Vertical line of frame S' at O'
    \coordinate (Op) at (2.75, 0);
    \draw (Op) -- (2.75, 2.35);
    \node[below=2pt] at (Op) {$O'$};
    \node[below=14pt] at (Op) {\small 观察者};

    % Frame labels
    \node at (1.35, 1.75) {$S$};
    \node at (3.45, 1.75) {$S'$};

    % Ray from source to observer
    \draw (src) -- (Op);

    % Angle theta at source
    \draw (src) ++(-90:0.52) arc (-90:{-90 + atan2(2.75, 2.0)}:0.52);
    \node at ($(src) + ({-90 + 0.45*atan2(2.75, 2.0)}:0.84)$) {$\theta$};

    % Distance d
    \draw (-0.45, 0) -- (-0.12, 0);
    \draw (-0.45, 2.0) -- (-0.12, 2.0);
    \node (dnode) at (-0.30, 1.0) {$d$};
    \draw[->] (dnode.north) -- (-0.30, 1.96);
    \draw[->] (dnode.south) -- (-0.30, 0.04);

    % Distance x
    \node (xnode) at (1.375, 0.55) {$x$};
    \draw[->] (xnode.west) -- (0.04, 0.55);
    \draw[->] (xnode.east) -- (2.71, 0.55);
\end{tikzpicture}

\end{center}
\end{problem}
\clearpage


\begin{problem}{12.2.31}
宇宙飞船上有一个发射器和一个接收器,它以恒定速度远离地球期间向地球发出一信号脉冲,飞船上的钟 40s 后收到被地球反射回来的信号,接收到的频率是发射频率的一半:

(1) 脉冲从地球反射时, 在飞船上测量地球位于何处?

(2) 飞船相对于地球的速度为何值?

(3) 脉冲被飞船接收时, 在地球上测量飞船位于何处?
\end{problem}
\clearpage


\begin{problem}{12.2.32}
在一惯性系中的观察者 O 观察到, 频率为 $\nu$ 的光子以入射角 $\theta_{i}$ 入射到相对于 O 以速度 v 在 x 方向运动的平面镜上, 求光子的反射角 $\theta_{r}$ 和反射光子的频率 $\nu_{r}$ 与 $\theta_{i}, \nu$ 的关系. 如果平面镜以速度 v 在 y 方向运动, 结果又将如何(图 12.16)?
\begin{center}
\begin{tikzpicture}[>=Stealth, line width=0.8pt]
    % Coordinate axes on the left
    \draw[->] (0, 0) -- (1.45, 0) node[below=2pt] {$x$};
    \draw[->] (0, 0) -- (0, 1.75) node[left=2pt] {$y$};

    % Moving vertical mirror with back hatching on the left
    \coordinate (M) at (2.45, 0.95);
    \draw (2.45, -0.55) -- (2.45, 2.25);
    \foreach \y in {-0.45, -0.25, ..., 2.15} {
        \draw (2.45, \y) -- (2.30, \y - 0.14);
    }

    % Normal to the mirror
    \draw (M) -- (4.85, 0.95);

    % Incident photon ray
    \draw[postaction={decorate}, decoration={markings, mark=at position 0.5 with {\arrow{Stealth}}}]
        (4.65, -0.22) -- (M);
    \node[below=3pt] at (3.85, 0.25) {$\nu$};

    % Reflected photon ray
    \draw[postaction={decorate}, decoration={markings, mark=at position 0.55 with {\arrow{Stealth}}}]
        (M) -- (4.55, 1.92);
    \node[above=3pt] at (3.95, 1.65) {$\nu_r$};

    % Incidence and reflection angles
    \draw (M) ++(-28:0.65) arc (-28:0:0.65);
    \node at ($(M) + (-14:1.00)$) {$\theta_i$};
    \draw (M) ++(0:0.75) arc (0:25:0.75);
    \node at ($(M) + (12.5:1.12)$) {$\theta_r$};

    % Mirror velocity v
    \draw[->] (2.45, -0.18) -- (3.30, -0.18) node[midway, below=2pt] {$v$};
\end{tikzpicture}

\end{center}
\end{problem}
\clearpage


\begin{problem}{12.2.33}
一面镜子在真空中以速度 v 沿 + x 方向运动, 一束频率为 $\nu$ 的光从 $x = +\infty$ 处垂直入射到镜面上, 求:

(1) 反射光的频率；

(2) 每个反射光子的能量；

(3) 若入射光的平均能流密度为 $I_{i}(\mathrm{W}/\mathrm{m}^{2})$ ，求反射光的平均能流密度.
\end{problem}
\clearpage


\begin{problem}{12.2.34}
在弗雷德·霍尔的一本小说的末尾的简化了的说法中,一个英雄以高的洛伦兹因子在与银河系平面成直角的方位飞行时,他说他似乎在一个蓝边红体的“金鱼缸”内部朝着缸口飞行.费曼打赌说,来自银河系的光看来不会是那样.你看谁正确?取 $\beta=0.99$ 和在银河参考系中观测边缘的角 $\theta$ 为 $45^{\circ}$ ,如图 12.17(a)所示.

(a)

(b)

(1) 导出或直接写出相对论性的光行差表达式, 并用它计算图 12.17(b) 中的 $\theta'$ , 即计算在宇宙飞船中看来自银河边缘的光的方向;

(2) 导出或直接写出相对论的多普勒效应公式, 并用它计算来自边缘的光的频率比 $\frac{\nu'}{\nu}$ ;

(3) 取足够多的角 $\theta$ 计算 $\theta'$ 和 $\frac{\nu'}{\nu}$ , 判定谁赌赢了.
\begin{center}
\begin{tikzpicture}[>=Stealth, line width=0.8pt]
    % Frame S' axes and velocity v
    \draw[->] (-1.15, 0) -- (-2.15, 0) node[midway, above=2pt] {$v$};
    \draw[->] (-1.15, 0) -- (2.75, 0) node[below=2pt] {$x'$};
    \draw[->] (0, 0) -- (0, 2.05) node[left=2pt] {$y'$};

    % Apparent Milky Way arc around the spacecraft
    \draw (45:1.15) arc (45:315:1.15);
    \node[above left=2pt] at (-0.95, 0.95) {\small 银河系};

    % Dashed line of sight and angle theta'
    \draw[dashed] (142:1.15) -- (-38:1.15);
    \draw (0:0.42) arc (0:142:0.42);
    \node at (58:0.65) {$\theta'$};

    % Pointer to spacecraft at origin
    \draw[->] (-0.65, -1.55) -- (-0.06, -0.14);
    \node[below=2pt] at (-0.65, -1.55) {\small 宇宙飞船};

    % Frame label
    \node at (1.85, 1.55) {\small $S'\text{（飞船参考系）}$};
\end{tikzpicture}

\end{center}
\begin{center}
\begin{tikzpicture}[>=Stealth, line width=0.8pt]
    % Frame S axes
    \draw[->] (0, 0) -- (3.85, 0) node[below=2pt] {$x$};
    \draw[->] (0, 0) -- (0, 2.25) node[right=2pt] {$y$};

    % Galactic plane ellipse
    \draw (0, 0) ellipse (0.55 and 1.55);
    \node[below=4pt] at (0, -1.58) {\small 银河系平面};

    % Spacecraft on x-axis moving at v
    \coordinate (ship) at (2.05, 0);
    \fill (ship) circle (2pt);
    \draw[->] (ship) -- (2.90, 0) node[midway, below=2pt] {$v$};

    % Pointer to spacecraft
    \draw[->] (2.55, -0.82) -- (2.10, -0.10);
    \node[below=2pt] at (2.55, -0.82) {\small 宇宙飞船};

    % Dashed line from edge of galactic plane through spacecraft
    \draw[dashed] (0.20, -1.44) -- (2.95, 0.70);
    \draw (ship) ++(0:0.48) arc (0:38:0.48);
    \node at ($(ship) + (19:0.75)$) {$\theta$};

    % Frame label
    \node at (2.15, 1.55) {\small $S\text{（银河参考系）}$};
\end{tikzpicture}

\end{center}
\end{problem}
\clearpage


\section{狭义相对论的动力学}

\begin{problem}{12.3.1}
两个静质量均为 $m$ 的粒子沿同一方向运动, 分别具有动量 $5mc$ 和 $10mc$ . 问

从较慢的粒子看,较快的粒子的速度多大?反之呢?
\end{problem}
\clearpage


\begin{problem}{12.3.2}
一静止质量为 m、荷电为 e 的静止带电粒子被场强为 E 的均匀电场加速了一段时间 t，获得了可与真空中光速相比的速度。问：

(1) 加速结束时, 粒子的动量和速度多大?

(2) 若此粒子不稳定, 在其静止参考系中以平均寿命 $\tau$ 衰变. 当此粒子以上述速度匀速运动时, 一个静止的观察者测到的粒子平均寿命多大?
\end{problem}
\clearpage


\begin{problem}{12.3.3}
(1) 一宇宙线质子与一静止质子碰撞形成一个高速运动 $\left[\gamma = \frac{1}{\sqrt{1 - \left(\frac{v}{c}\right)^2}} = 1000\right]$ 的激发系统. 在此系统中, 介子以速率 $\beta' c\left(\beta' = \frac{v'}{c}\right)$ 被射出. 如在动参考系中,介子射出方向与v的夹角为 $\theta'$ ,求在实验室参考系中观察, $\theta$ 多大?

(2) 把(1)问中得到的结果应用于在动系中, 以动量 $0.5\mathrm{GeV} / \mathrm{c}$ 被射出的介子 (静能为 $140\mathrm{MeV}$ ), 若 $\theta' = 90^{\circ}, \theta$ 多大? 在实验室中观察到的 $\theta_{\max}$ 多大?
\end{problem}
\clearpage


\begin{problem}{12.3.6}
静止质量 $m$ 、荷电 $e$ 的粒子以相对论速率 $v$ 在半径为 $R$ 的圆轨道上运动，轨道平面垂直于均匀静磁场 $\bm{B}$ ，如图 12.19(a) 所示，以恒定速率 $v$ 沿 $-y$ 轴运动中的观察者 $O'$ 看到的轨道如图 12.19(b) 所示，在两图中 $a, b, c, d, e$ 诸点是相应的。图中 $\bm{B}$ 的画法，表明 $\bm{B}$ 沿 $+z$ 方向。

(a)

(b)

(1) 由 $O'$ 测得距离 $y_d' - y_b'$ 多大?

(2) 证明 $c$ 点是 $S'$ 系中的静止点, 在此点, $\frac{\mathrm{d}^2 x'}{\mathrm{d} t'^2}$ 、 $\frac{\mathrm{d}^2 y'}{\mathrm{d} t'^2}$ 多大?

(3) 由 $O'$ 看来, 是什么引起在 $c$ 点处的加速度的?
\begin{center}
\begin{tikzpicture}[>=Stealth, line width=0.8pt]
    % Axes in S' (y' horizontal right, x' vertical down)
    \draw[->] (-1.35, 0) -- (4.15, 0) node[below=2pt] {$y'$};
    \draw[->] (0, 1.25) -- (0, -1.45) node[right=2pt] {$x'$};
    \node[below left=2pt] at (0, 0) {$O'$};

    % Cycloid-like trajectory with cusp at c (phi = pi)
    \draw[domain=-1.95:3.14159, samples=90, smooth]
        plot ({0.48*(\x + sin(\x r))}, {-0.78*cos(\x r)});
    \draw[domain=3.14159:6.85, samples=90, smooth]
        plot ({0.48*(\x + sin(\x r))}, {-0.78*cos(\x r)});

    % Marked points a, b, c, d, e
    \coordinate (pa) at (0, -0.78);
    \coordinate (pb) at (1.234, 0);
    \coordinate (pc) at (1.508, 0.78);
    \coordinate (pd) at (1.782, 0);
    \coordinate (pe) at (3.016, -0.78);

    \foreach \p in {pa, pb, pc, pd, pe} {
        \fill (\p) circle (2pt);
    }

    \node[below right=1pt] at (pa) {$a$};
    \node[above left=1pt] at (pb) {$b$};
    \node[above=3pt] at (pc) {$c$};
    \node[above right=1pt] at (pd) {$d$};
    \node[below right=1pt] at (pe) {$e$};
\end{tikzpicture}

\end{center}
\begin{center}
\begin{tikzpicture}[>=Stealth, line width=0.8pt]
    % Axes in S (y horizontal right, x vertical down)
    \draw[->] (-1.85, 0) -- (1.95, 0) node[below=2pt] {$y$};
    \draw[->] (0, 1.85) -- (0, -1.95) node[right=2pt] {$x$};
    \node[below left=2pt] at (0, 0) {$O$};

    % Circular orbit of radius R, counterclockwise (a -> b -> c -> d -> e)
    \draw[postaction={decorate}, decoration={markings, mark=at position 0.875 with {\arrow{Stealth}}}]
        (0, 0) circle (1.25);
    \node[below right=2pt] at (0.88, -0.88) {$v$};

    % Radius vector R
    \draw[->] (0, 0) -- (32:1.25);
    \node[above left=1pt] at (32:0.68) {$R$};

    % Marked points a, b, c, d, e
    \fill (0, -1.25) circle (2pt) node[below right=1pt] {$ae$};
    \fill (1.25, 0) circle (2pt) node[above right=1pt] {$b$};
    \fill (0, 1.25) circle (2pt) node[above right=1pt] {$c$};
    \fill (-1.25, 0) circle (2pt) node[below left=1pt] {$d$};

    % Magnetic field B out of page (+z)
    \draw (1.35, 1.55) circle (0.16);
    \fill (1.35, 1.55) circle (1.8pt);
    \node[right=3pt] at (1.51, 1.55) {$\pmb{B}$};
\end{tikzpicture}

\end{center}
\end{problem}
\clearpage


\begin{problem}{12.3.7}
试计算相对论粒子的“纵向质量”和“横向质量”. “纵向质量”是指作用力与运动方向平行时力与粒子的加速度之比, “横向质量”是指力与运动方向垂直时力与加速度之比.
\end{problem}
\clearpage


\begin{problem}{12.3.8}
一静止质量为 $m$ 的粒子最初沿 $x$ 轴运动, 速度为 $v_0, t = 0$ 后受到一个沿 $y$ 方向的恒力 $F$ 作用, 求它在任意时刻 $t$ 的速度, 并证明 $t \to \infty$ 时, $|v| \to c$ .
\end{problem}
\clearpage


\begin{problem}{12.3.9}
一静止质量为 $m$ 、荷电为 $e$ 的电子在与均匀磁场垂直的平面内运动, 如果忽略辐射, 电子的轨道为半径为 $R$ 的圆. 设电子能量 $E \gg mc^2$ , 求磁感应强度 $B$ 的表达式. 如 $R = 30\mathrm{m}$ , $E = 2.5 \times 10^{9}\mathrm{eV}$ . 求出 $B$ 的数值.
\end{problem}
\clearpage


\begin{problem}{12.3.10}
（1）静止的 $\mu$ 子平均寿命 $T_{0}$ 为 $10^{-6}s$ ，静止质量为 $100MeV/c^{2}$ ，大气中高度 h 约为 $10^{4}m$ 处的 $\mu$ 子要到达地面，它需要带有多大能量 E？

(2) 假定地球具有 $1 \times 10^{-4} \mathrm{T}$ 的磁场, 方向沿地轴方向, 延伸范围在地面以上 $10^{2} \mathrm{~m}$ 范围内, 问在赤道上空垂直入射能量为 (1) 中算出的 $E$ 值的 $\mu$ 子被磁场偏转的大小和方向 (可不计万有引力作用).
\end{problem}
\clearpage


\begin{problem}{12.3.11}
导出自由的相对论粒子的哈密顿函数和拉格朗日函数.
\end{problem}
\clearpage


\begin{problem}{12.3.12}
导出具有动量 $p = \frac{m_0v}{\sqrt{1 - \frac{v^2}{c^2}}}$ 的粒子处在电磁场

\[
\bm {E} = - \nabla \phi - \frac {\partial \bm {A}}{\partial t}, \quad \bm {B} = \nabla \times \bm {A}
\]

中时的哈密顿函数.
\end{problem}
\clearpage


\begin{problem}{12.3.13}
一个静止质量为 $m$ 、荷电为 $e$ 、速度为 $\bm{v}$ 的相对论粒子，在一个矢势为 $\bm{A}$ 的电磁场中运动，其拉格朗日函数为

\[
L = - m c ^ {2} \sqrt {1 - \frac {v ^ {2}}{c ^ {2}}} + e A \cdot v
\]

对于位于原点的、沿 $z$ 轴的磁矩为 $M_0$ 的磁偶极子产生的磁场可用矢势 $\mathbf{A} = \frac{\mu_0}{4\pi}\frac{M_0\sin\theta}{r^2}\mathbf{e}_{\varphi}$ 描写，其中 $r, \theta, \varphi$ 为球坐标.

(1) 求 $p_{\varphi}$ ，并证明它是一个运动常数；

(2) 如果上述矢势被 $A' = A + \nabla \chi(r, \theta, \varphi)$ 代替， $\chi$ 为一任意标量函数， $p_{\varphi}$ 将如何变化，是否还是运动常数？
\end{problem}
\clearpage


\begin{problem}{12.3.14}
证明: 如果恒定磁场是平面场: $A_{x}=0, A_{y}=0, A_{z}=A_{z}(x, y)$ ，则带电粒子在该场中运动时，

\[
\frac {m v _ {z}}{\sqrt {1 - \beta^ {2}}} + q A _ {z} = \text {常量}
\]
\end{problem}
\clearpage


\begin{problem}{12.3.15}
证明: 如果恒定磁场具有轴对称性, 即 $A_{r} = A_{z} = 0, A_{\varphi} = A_{\varphi}(r, z)$ , 则带电粒子在该场中运动时.

\[
\frac {m r ^ {2} \dot {\varphi}}{\sqrt {1 - \beta^ {2}}} + q r A _ {\varphi} = \text {常量}
\]
\end{problem}
\clearpage


\begin{problem}{12.3.16}
众所周知,行星都围绕太阳在椭圆轨道上运动,但如果仅考虑狭义相对论的效应,轨道是进动的椭圆.

\[
\frac {1}{r} = \frac {1}{r _ {0}} \{1 + \varepsilon \cos [ \alpha (\varphi - \varphi_ {0}) ] \}
\]

$\alpha\neq1(\alpha=1$ 相应于无进动的经典的结果).

(1) 导出这一方程, 用轨道的基本常数 (如能量、角动量等) 表示 $\alpha$ 和 $r_0$ ;

(2) 已知水星的轨道的平均半径为 $58 \times 10^{6} ~km$ ，轨道周期为 88 天，计算每 100 年水星轨道进动的角度（当然，这个效应不能说明水星总的进动速率）.
\end{problem}
\clearpage


\begin{problem}{12.3.17}
试求在静止原子核 $Ze$ 的库仑场中, 静止质量为 $m$ 、带电 $-e$ 的电子的椭圆轨道的相对论修正.
\end{problem}
\clearpage


\begin{problem}{12.3.19}
若一粒子的动能等于它的静能,它的速度多大?
\end{problem}
\clearpage


\begin{problem}{12.3.20}
试用相对论粒子的动能 T 表示它的动量的大小 p，用相对论粒子的动量 p 表示它的速度 v。
\end{problem}
\clearpage


\begin{problem}{12.3.21}
一个静止质量为 $m_0$ 的相对论粒子具有能量 $E$ , 求该粒子的速率 $v$ , 讨论非相对论和极相对论两种极限情况.
\end{problem}
\clearpage


\begin{problem}{12.3.22}
(1) 计算与具有动量 $400\mathrm{GeV} / c$ 的质子同样速度的 $\pi$ 介子的动量. 在费米实验室, 动量为 $400\mathrm{GeV} / c$ 的质子打击靶时, 是产生 $\pi$ 介子的最可几动量, $\pi$ 介子的静止质量为 $0.14\mathrm{GeV} / c^2$ , 质子的静止质量为 $0.94\mathrm{GeV} / c^2$ ;

(2) $\pi$ 介子通过 $400\mathrm{m}$ 长的衰变管时, 一些 $\pi$ 介子衰变产生中微子束, 中微子探测器位于 $1\mathrm{km}$ 以外, 如图 12.20 所示. $\pi$ 介子的固有的平均寿命为 $2.6 \times 10^{-8}\mathrm{s}$ , 在 $400\mathrm{m}$ 长的距离上, $\pi$ 介子发生衰变占多大分数?

(3) 在 $\pi$ 介子静止参考系中的观察者测量, 衰变管的长度是多少?

(4) $\pi$ 介子衰变成一个 $\mu$ 子和一个中微子 $\nu$ (静止质量为零). 用相对论的能量动量关

系证明，在 $\pi$ 介子静止的参考系中，衰变产物的动量的大小 $q$ 由下式确定：

\[
\frac {q}{c} = \frac {M ^ {2} - m ^ {2}}{2 M}
\]

其中 M 是 $\pi$ 介子静止质量, m 是 $\mu$ 子的静止质量;

(5) 平均来讲, 中微子探测器离 $\pi$ 介子衰变点约 $1.2\mathrm{km}$ , 为了在 $\pi$ 介子静止参考系中前半球中产生的所有中微子都被探测到, 探测器的横截面的半径应多大?
\begin{center}
\begin{tikzpicture}[>=Stealth, line width=0.8pt]
    % Proton beam
    \draw[->] (0, 0) -- (0.68, 0);
    \node at (0.34, -0.35) {\scriptsize 质子束};
    \node at (0.34, -0.72) {\scriptsize $400\,\mathrm{GeV}$};

    % Target
    \fill (0.82, -0.15) rectangle (1.22, 0.15);
    \node at (1.02, 0.38) {\scriptsize 靶};

    % Pion decay pipe (400 m)
    \draw (1.42, 0.22) -- (3.12, 0.22);
    \draw (1.42, -0.22) -- (3.12, -0.22);
    \node at (2.27, 0) {\scriptsize $\pi$介子衰变管};

    % Focusing system pointer
    \draw[->] (1.52, -0.82) -- (1.52, -0.26);
    \node[below=2pt] at (1.52, -0.82) {\scriptsize 聚焦系统};

    % Shielding with hatched border and clear center
    \fill[pattern=north east lines] (3.12, -0.48) rectangle (5.28, 0.48);
    \draw (3.12, -0.48) rectangle (5.28, 0.48);
    \fill[white] (3.26, -0.35) rectangle (5.14, 0.35);
    \node at (4.20, 0) {\scriptsize 屏蔽罩};

    % Neutrino detector
    \draw (5.75, 0) circle (0.30);
    \draw[->] (5.75, -0.82) -- (5.75, -0.34);
    \node[below=2pt] at (5.75, -0.82) {\scriptsize 中微子探测器};

    % Top length dimensions (400 m and 1000 m)
    \draw (1.42, 0.58) -- (1.42, 1.12);
    \draw (3.12, 0.58) -- (3.12, 1.12);
    \draw (5.75, 0.58) -- (5.75, 1.12);

    \node (d1) at (2.27, 0.90) {\scriptsize $400\,\mathrm{m}$};
    \draw[->] (d1.west) -- (1.45, 0.90);
    \draw[->] (d1.east) -- (3.09, 0.90);

    \node (d2) at (4.435, 0.90) {\scriptsize $1000\,\mathrm{m}$};
    \draw[->] (d2.west) -- (3.15, 0.90);
    \draw[->] (d2.east) -- (5.72, 0.90);
\end{tikzpicture}

\end{center}
\end{problem}
\clearpage


\begin{problem}{12.3.23}
如果将 100kg 铜的温度升高 100℃, 它的质量增加了多少? (铜的比热容 c = 93cal/(kg·K))
\end{problem}
\clearpage


\begin{problem}{12.3.24}
在海拔 100km 的地球大气层中产生了一个静能为 140MeV 的 $\pi^{+}$ 介子，这个 $\pi^{+}$ 介子的总能量 $E=1.5\times10^{5}MeV$ ，竖直向下运动，按它自身参考系中测定，它在产生后 $2\times10^{-8}s$ 衰变，问它在海平面以上多大高度处发生衰变的？
\end{problem}
\clearpage


\begin{problem}{12.3.25}
一个假想的手电筒能将它的相当一部分静止质量转变成光, 发射相当准直的光束, 如果手电筒开始处于静止, 静止质量为 $m_0$ . 接通开关并允许它沿一直线自由地运动. 求当它相对于原静参考系具有速度 $v$ 时的静止质量 $m$ (不假定 $c \gg v$ ).
\end{problem}
\clearpage


\begin{problem}{12.3.26}
一个能量为 $E_{i}$ 的光子被一个静止质量为 $m_{\mathrm{e}}$ 的静止电子所散射，设散射后光子的能量为 $E_{f}$ ，入射光子与散射光子之间的夹角为 $\theta$ ，找出 $E_{i}, E_{f}$ 和 $\theta$ 之间的关系。
\end{problem}
\clearpage


\begin{problem}{12.3.27}
一个静止的 $\pi^{+}$ 介子衰变为 $\mu^{+}$ 子和中微子，三者的静止质量分别为 $m_{\pi 0}$ 、 $m_{\mu 0}$ 和零，求 $\mu^{+}$ 子和中微子的动能 $T_{\mu}, T_{\nu}$ .
\end{problem}
\clearpage


\begin{problem}{12.3.28}
静止质量均为 $9.11 \times 10^{-31} ~kg$ 的一个电子和一个正电子相碰撞，两粒子消失产生电磁辐射，若湮没前均静止，求辐射出的能量。
\end{problem}
\clearpage


\begin{problem}{12.3.29}
两个静止质量均为 $m_0$ 的粒子, 一个处于静止状态, 另一个动能为其静能的 6 倍, 两粒子做完全非弹性碰撞, 求复合粒子的静止质量和速度.
\end{problem}
\clearpage


\begin{problem}{12.3.30}
一个静止质量为 $m_0$ 、速率为 $0.8c$ 沿 $x$ 轴向正方向运动的粒子，衰变成两个静止质量均为 $m_0$ 的粒子，其中一个以 $0.6c$ 的速率沿 $y$ 轴向负方向运动。求：

(1) 另一个粒子的速率和运动方向；

(2) $m_{0}/M_{0}$ .
\end{problem}
\clearpage


\begin{problem}{12.3.31}
一总能量为 $E_0$ 、静止质量为 $m_0$ 的粒子，与一个具有相同静止质量的静止粒子发生完全弹性碰撞，散射后两粒子动能相等。求散射后两粒子速度的夹角，并讨论非相对论和极相对论两种极限情况。
\end{problem}
\clearpage


\begin{problem}{12.3.32}
一个能量为 $E_0$ 的反质子与一处于静止的质子相互作用, 产生两静止质量均为 $m_x$ 的粒子, 其中一个在实验室中在与入射方向成 $90^\circ$ 角的方向上被探测到, 求这个粒子的总能量 $E_s$ , 并证明它与 $m_x, E_0$ 均无关.
\end{problem}
\clearpage


\begin{problem}{12.3.33}
宇航员开亮一个普通的手电筒,将它扔入太空(电筒围绕它的轴稳定地自转),问两小时后电池用完之时,这个“光子火箭”获得多大的附加速度?
\end{problem}
\clearpage


\begin{problem}{12.3.34}
等效原理断言引力质量和惯性质量相等. 光子有不等于零的引力质量吗? 设光子朝着地球落下距离 10m, 计算对光子频率的影响. 用什么实验技术可以测出这个频率变化?
\end{problem}
\clearpage


\begin{problem}{12.3.35}
考虑一个能量非常高的散射实验, 参与散射的两个粒子具有相同的静止质量 $m_{0}$ , 其中一个原来处于静止, 另一个以动量 p 入射:

(1) 求质心的速度 $v^{*}$ ;

(2) 在极相对论极限 $pc \gg m_{0}c^{2}$ 下，求此系统在质心平动参考系中的总能量 $E^{*}$ .
\end{problem}
\clearpage


\begin{problem}{12.3.36}
(1) 静止质量 $m_{1} = 1\mathrm{g}$ 以0.9倍光速运动的粒子，与静止质量 $m_{2} = 10\mathrm{g}$ 的静止粒子发生正碰，并嵌入其中，产生的复合粒子的静止质量和速度多大？

(2) 假定 $m_{1}$ 静止, $m_{2}$ 应以多大的速度运动, 才能产生与(1)有相同静止质量的复合粒子?

(3) 若 $m_{1}$ 静止, $m_{2}$ 应以多大的速度运动, 才能产生与(1)有相同速度的复合粒子.
\end{problem}
\clearpage


\begin{problem}{12.3.37}
静能为 494MeV 的 K 介子衰变成静能为 106MeV 的 $\mu$ 子和静能为零的中微子，求由静止的 K 介子衰变成的 $\mu$ 子和中微子的动能.
\end{problem}
\clearpage


\begin{problem}{12.3.38}
可以通过辐射 $\gamma$ 量子的方法使被激励的原子核返回基态. 设基态的质量即原子核的静止质量为 $m_{0}$ ，激发能(处于该激发态时原子核的能量减去处于基态时的能量)为 $\Delta E$ ，求辐射出的 $\gamma$ 量子的频率.
\end{problem}
\clearpage


\begin{problem}{12.3.39}
一个静止质量为 $M_0$ 的静止粒子裂变为静止质量分别为 $m_{10}$ 和 $m_{20}$ 的两粒子，裂变能 $M_0c^2 - (m_{10} + m_{20})c^2$ 在这两个粒子间如何分配？
\end{problem}
\clearpage


\begin{problem}{12.3.40}
静止的电子偶湮没时产生两个光子,如果其中一个光子再与另一个静止电子发生碰撞,它能给予这个电子的最大速度多大?
\end{problem}
\clearpage


\begin{problem}{12.3.41}
在气泡室中经常能观察到由光子产生电子-正电子对的反应. 试证明除非有其他粒子的介入, 比如一个核子, 则这样的转化是不可能的.
\end{problem}
\clearpage


\begin{problem}{12.3.42}
一电子束被固定的散射靶散射,如图 12.24 所示,散射是弹性的,每个电子具有能量 $E=\frac{5}{3}m_{0}c^{2}(m_{0}$ 为电子的静止质量),电子束流量为每秒钟流过 Q 个电子.求:

(1) 入射电子的速度；

(2) 电子作用在散射靶上的力的大小和方向.
\begin{center}
\begin{tikzpicture}[>=Stealth, thick, line cap=round, line join=round]
  \coordinate (O) at (0, 0);
  \coordinate (N) at (-3.8, 0);
  \coordinate (A) at (-2.8, 2.8);
  \coordinate (B) at (-2.8, -2.8);

  % 散射靶（右侧竖直墙面及斜线阴影）
  \draw[very thick] (0, -2.4) -- (0, 2.4);
  \foreach \y in {-2.2, -1.85, ..., 2.15} {
    \draw[thin] (0, \y) -- ++(0.25, 0.25);
  }

  % 法线（水平虚线）
  \draw[dashed] (N) -- (O);

  % 入射与散射电子束
  \draw[postaction={decorate}, decoration={markings, mark=at position 0.52 with {\arrow{Stealth}}}]
    (A) -- (O);
  \draw[postaction={decorate}, decoration={markings, mark=at position 0.55 with {\arrow{Stealth}}}]
    (O) -- (B);

  % 45 度角标注
  \pic[draw, thin, angle radius=1.1cm, angle eccentricity=1.4, "$45^\circ$"] {angle=A--O--N};
  \pic[draw, thin, angle radius=1.1cm, angle eccentricity=1.4, "$45^\circ$"] {angle=N--O--B};
\end{tikzpicture}

\end{center}
\end{problem}
\clearpage


\begin{problem}{12.3.43}
一个静止质量为 m、能量为 E 的粒子，跟另一个静止质量为 M 的静止粒子发生弹性碰撞，求碰撞后这个粒子的能量 $E_{1}$ 与散射角 $\theta$ 的关系，讨论 m=0 的情况（康普顿效应）.
\end{problem}
\clearpage


\begin{problem}{12.3.44}
以速度 $v$ 运动、静止质量为 $m$ 的 $\pi^0$ 介子裂变成两个 $\gamma$ 量子，设在 $\pi^0$ 介子静止的参考系内， $\gamma$ 量子按飞散方向的分布是各向同性的，试确定相对于实验室参考系的下列各量：

(1) 其中一个 $\gamma$ 量子按与 $\pi^{0}$ 介子运动方向成 $\theta$ 角飞散的概率；

(2) 一个 $\gamma$ 量子以 $\theta$ 角飞散, 另一个 $\gamma$ 量子飞散的方向;

(3) 按(2)飞散的两个 $\gamma$ 量子的能量.
\end{problem}
\clearpage


\begin{problem}{12.3.45}
一个能量为 $E_{\nu}$ 的光子与一个能量为 $E \gg m_0 c^2 (m_0$ 为静止质量)的电子发生碰撞.

(1) 若在实验室参考系中光子运动方向与电子运动方向的夹角为 $\theta$ ，碰撞前光子在电子参考系中的能量 $E_{v}^{\prime}$ 等于多少？

(2) 若 $E_{\nu}^{\prime} \ll m_{0}c^{2}$ ，电子的反冲可以略去，光子在电子参考系中的能量不会因为碰撞而改变，问碰撞后光子在实验室参考系中能量的最小值和最大值各是多少？
\end{problem}
\clearpage


\begin{problem}{12.3.46}
相距为 $a$ 的两个点电荷 $q_{1}$ 和 $q_{2}$ , 在垂直于它们的连线的方向上以匀速度 $\bm{v}$ 运动, 试确定它们的相互作用力.
\end{problem}
\clearpage


\section{四维矢量}

\begin{problem}{12.4.1}
(1) 已知 $(\bm{r},ct)$ 是一个相对论性四维矢量, 证明 $(\bm{k},\omega)$ 是相对论性四维矢量.

(2) 已知一原子在静止时放射出角频率为 $\omega_0$ 的光, 且这原子以速率 $v$ 朝向或离开一观察者运动, 利用洛伦兹变换就两种情况 (朝向或离开) 导出观察者观察到的角频率公式.
\end{problem}
\clearpage


\begin{problem}{12.4.2}
写出四维位置矢量的洛伦兹变换,导出四维动量矢量的洛伦兹变换.
\end{problem}
\clearpage


\begin{problem}{12.4.3}
试证对于两个粒子的系统, 若 p、E 分别为系统的总动量和总能量, 则

\[
p _ {x} ^ {2} + p _ {y} ^ {2} + p _ {z} ^ {2} - \left(\frac {E}{c}\right) ^ {2}
\]

为洛伦兹变换下的不变量.
\end{problem}
\clearpage


\begin{problem}{12.4.4}
问两束动能均为 $T$ 的质子迎头碰撞的反应与多大动能的单束质子跟静止的质子碰撞时的反应相同.
\end{problem}
\clearpage


\begin{problem}{12.4.5}
一个动量为 p 的光子与一个静止质量为 m 的静止粒子发生碰撞.

(1) 在动量中心参考系中, 光子和粒子的总能量多大?

(2) 在动量中心系中, 粒子动量的量值多大?

(3) 如果光子向后弹性散射, 末态光子在实验室参考系中的动量多大?
\end{problem}
\clearpage


\begin{problem}{12.4.6}
一光子与静止的靶质子发生碰撞产生如下的反应.

\[
\gamma + p \rightarrow p + \psi^ {\prime}
\]

近似用 $\psi'$ 粒子的静止质量为 $4M_{p}, M_{p}$ 是质子的静止质量，求在实验室参考系中：

(1) 产生上述反应的光子所必须具有的最小能量 $E_0$ (可以 $M_{\mathrm{p}}c^{2}$ 为单位);

(2) 当光子能量 E 恰是上述阈能 $E_{0}$ 时产生的 $\psi'$ 粒子的速度.

(1) 在光子能量恰为阈能 $E_0$ 时, 其动量 $p = \frac{E_0}{c}$ , 与静止质子发生碰撞产生 $\psi'$ 粒子的反应时, 在动量中心参考系中, 质子和 $\psi'$ 粒子都是静止的 (在实验室参考系中, 由于要遵从动量守恒定律不可能都静止).

由 $E^{2}-c^{2}p^{2}$ 是洛伦兹变换下的不变量以及孤立系统的守恒量，反应前考虑实验室参考系, 这个守恒量为 $(E_0 + M_{\mathrm{p}}c^2)^2 - c^2 p^2 = (E_0 + M_{\mathrm{p}}c^2)^2 - E_0^2$ (用了系统的动量 $p$ 等于光子的动量 $\frac{E_0}{c}$ ).

反应后考虑动量中心系，这个守恒量为 $(4M_{\mathrm{p}}c^{2} + M_{\mathrm{p}}c^{2})^{2} = 25(M_{\mathrm{p}}c^{2})^{2}$ 

\[
(E _ {0} + M _ {\mathrm{p}} c ^ {2}) ^ {2} - E _ {0} ^ {2} = 2 5 (M _ {\mathrm{p}} c ^ {2}) ^ {2}
\]

解出

\[
E _ {0} = 1 2 M _ {\mathrm{p}} \mathrm{c} ^ {2}
\]

说明: 在动量中心系中, 系统的动量为零, 如果两粒子不是静止的, 显然在质心系中, 上述守恒量将大于 $25(M_{\mathrm{p}}c^{2})^{2}$ , 入射光子的能量 E 将大于 $E_{0}$ , 故 $E_{0}$ 确是阈能.

(2) 入射光子具有阈能 $E_{0}$ 时, 产生的 $\psi'$ 粒子在动量中心系中速度为零, 用速度变换公式可知, 在实验室参考系中, 它的速度就是动量中心参考系的速度.

在动量中心系中,系统的动量 $p' = 0$ , 在实验室参考系中, 系统的动量为碰撞前光子的动量 $p = \frac{E_{0}}{c}$ , 由四维动量的洛伦兹变换

\[
\begin{array}{r l} 0 & = p ^ {\prime} = \gamma \left(p - \frac {\beta E}{c}\right) \\ \beta & = \frac {p c}{E} = \frac {E _ {0}}{E _ {0} + M _ {\mathrm{p}} c ^ {2}} = \frac {1 2 M _ {\mathrm{p}} c ^ {2}}{1 3 M _ {\mathrm{p}} c ^ {2}} = \frac {1 2}{1 3} \\ v & = \beta c = \frac {1 2}{1 3} c \end{array}
\]
\end{problem}
\clearpage


\begin{problem}{12.4.7}
一个静止质量为 $m$ 的粒子与一个静止质量为 $m_{1}$ 的静止粒子碰撞, 发生有很多粒子产生的反应, 这些粒子的总静止质量为 $M$ . 如果 $m + m_{1} < M$ , 那么, 当入射粒子的动能不够大时, 反应便不能进行. 试求出反应的阈动能 $T_{0}$ , 即反应刚能发生入射粒子应具有的动能.
\end{problem}
\clearpage


\begin{problem}{12.4.8}
一个入射核子与另一个静止的核子碰撞产生 $\pi$ 介子, 反应为

\[
\mathrm{N} + \mathrm{N} \rightarrow \mathrm{N} + \mathrm{N} + \pi
\]

求此反应的阈动能 $T_{0}$
\end{problem}
\clearpage


\begin{problem}{12.4.9}
试证明: 只有当静止质量 $m_{1} \neq 0$ 的粒子参加反应(尽管粒子在反应后仍然存在)时, 一个 $\gamma$ 量子才有可能转化为一个电子和正电子偶, 试求这反应能发生的阈能 $E_{0}$ , 这个粒子在反应前后没发生任何变化吗?
\end{problem}
\clearpage


\begin{problem}{12.4.10}
在高能质子-质子碰撞中,一个或两个质子可以分解成一个质子和几个带电 $\pi$ 介子,其反应为:

(a) $\mathrm{p} + \mathrm{p}\rightarrow \mathrm{p} + (\mathrm{p} + n\pi)$ 

(b) $\mathrm{p} + \mathrm{p}\rightarrow (\mathrm{p} + n\pi) + (\mathrm{p} + m\pi).$ 

这里 $n$ 和 $m$ 是产生 $\pi$ 介子的数目.

在实验室参考系中, 一个总能量为 E 的入射质子(下文称“入射粒子”)打击一个静止的质子(下文称“靶”), 求下列各情形入射粒子所需的最小能量.

(1) 发生反应(a)，靶分解为一个质子和 4 个 $\pi$ 介子；

(2) 发生反应(a)，入射粒子分解为一个质子和 4 个 $\pi$ 介子；

(3) 发生反应(b)，两个质子都分解成一个质子和 4 个 $\pi$ 介子.

\[
m _ {\pi} c ^ {2} = 0. 1 4 0 \mathrm{GeV}, \quad m _ {\mathrm{p}} c ^ {2} = 0. 9 3 8 \mathrm{GeV}
\]
\end{problem}
\clearpage


\begin{problem}{12.4.11}
考虑 $\pi$ 介子的光致反应

\[
\gamma + \mathrm{p} \rightarrow \mathrm{p} + \pi^ {0}
\]

质子的静能为 938MeV, $\pi^{0}$ 介子的静能为 135MeV.

(1) 如果初始质子相对实验室静止, 求出此反应要能进行 $\gamma$ 射线在实验室参考系中的阈能;

(2) 各向同性的 3K 字宙黑体辐射的平均光子能量约为 $10^{-3}\mathrm{eV}$ . 考虑一个质子和一个能量为 $10^{-3}\mathrm{eV}$ 的光子间的正碰, 求出此 $\pi$ 介子光致反应能进行的最小质子能量.
\end{problem}
\clearpage


\begin{problem}{12.4.12}
一动量为 $5m_{\pi}c$ 的 $\pi$ 介子与开始处于静止的质子 $(m_{\mathrm{p}} = 7m_{\pi})$ 发生弹性碰撞，求：

(1) 动量中心参考系的速度；

(2) 在动量中心系中 $\pi$ 介子和质子的总能量；

(3) 在动量中心系中入射 $\pi$ 介子的动量.
\end{problem}
\clearpage


\begin{problem}{12.4.13}
粒子物理中特别令人感兴趣的是高能情况下的弱相互作用, 让 $\pi$ 介子和 K 介子在飞行中衰变产生中微子束. 假定一个动量为 $200\mathrm{GeV} / c$ 的 $\pi$ 介子通过衰变 $\pi \rightarrow \mu + \nu$ 产生中微子, $\pi$ 介子在 $\pi$ 介子静止参考系中的寿命为 $\tau_{\pi \pm} = 2.60 \times 10^{-8}\mathrm{s}$ , 其静能为 $139.6\mathrm{MeV}, \mu$ 子的静能为 $105.7\mathrm{MeV}$ , 中微子是无静止质量的.

(1) 计算 $\pi$ 介子衰变前在实验室中行走的平均距离；

(2) 计算 $\mu$ 子在实验室中的最大角度(相对于 $\pi$ 介子的运动方向);

(3) 计算中微子在实验室参考系中能够具有的最小和最大动量.
\end{problem}
\clearpage


\begin{problem}{12.4.14}
粒子物理中,粒子间的相互作用被设想为如图 12.27 所示,经一个粒子的交换而引起的,证明这样交换粒子不是真实的,而是虚构的.
\begin{center}
\begin{tikzpicture}[>=Stealth, thick, line cap=round, line join=round,
  fermion/.style={postaction={decorate}, decoration={markings, mark=at position 0.52 with {\arrow{Stealth}}}}]
  \coordinate (V1) at (-1.4, 0);
  \coordinate (V2) at (1.4, 0);

  % 相互作用顶点
  \fill (V1) circle (2.2pt);
  \fill (V2) circle (2.2pt);

  % 左侧粒子世界线
  \draw[fermion] (-2.3, -1.25) -- (V1);
  \draw[fermion] (V1) -- (-2.6, 1.1);

  % 右侧粒子世界线
  \draw[fermion] (2.8, -0.55) -- (V2);
  \draw[fermion] (V2) -- (2.15, 1.2);

  % 交换粒子（波浪线）
  \draw[decorate, decoration={snake, amplitude=2.0mm, segment length=8.5mm}] (V1) -- (V2);
\end{tikzpicture}

\end{center}
\begin{center}
\begin{tikzpicture}[>=Stealth, thick, line cap=round, line join=round,
  fermion/.style={postaction={decorate}, decoration={markings, mark=at position 0.52 with {\arrow{Stealth}}}}]
  \coordinate (V1) at (-1.4, 0);
  \coordinate (V2) at (1.4, 0);

  % 相互作用顶点及编号
  \fill (V1) circle (2.2pt) node[below=3pt] {$1$};
  \fill (V2) circle (2.2pt) node[below=3pt] {$2$};

  % 左侧粒子初末态
  \draw[fermion] (-2.3, -1.25) -- node[pos=0.45, below right=0pt and 3pt] {$\boldsymbol{p}_1$} (V1);
  \draw[fermion] (V1) -- node[pos=0.62, above right=0pt and 3pt] {$\boldsymbol{p}_1'$} (-2.6, 1.1);

  % 右侧粒子初末态
  \draw[fermion] (2.8, -0.55) -- node[pos=0.45, below left=2pt and 1pt] {$\boldsymbol{p}_2$} (V2);
  \draw[fermion] (V2) -- node[pos=0.62, above left=0pt and 3pt] {$\boldsymbol{p}_2'$} (2.15, 1.2);

  % 交换粒子（波浪线，箭头由 2 指向 1）
  \draw[decorate, decoration={snake, amplitude=2.0mm, segment length=8.5mm},
        postaction={decorate, decoration={markings, mark=at position 0.50 with {\arrow{Stealth[reversed]}}}}]
    (V1) -- (V2);
\end{tikzpicture}

\end{center}
\end{problem}
\clearpage


\begin{problem}{12.4.15}
在一个相对论性的核-核碰撞的简化模型中, 一个静止质量为 $m_{1}$ 、速率为 $\beta_{1}c$ 的核与一个静止质量为 $m_{2}$ 的静止的靶核做正碰, 这个系统的动量中心参考系的速度为 $\beta_{0}c$ , 在此参考系中系统的能量为 $E_{0}$ .

(1) 用相对论导出 $\beta_{0}$ 和 $E_{0}$ 的表达式；

(2) 就一个 $^{40}Ar$ 核以 $\beta_{1}=0.8$ 与一个静止的 $^{238}U$ 核碰撞, 计算 $\beta_{0}$ 和 $E_{0}$ (以 MeV 为单位);

(3) 在 $\mathrm{Ar} + \mathrm{U}$ 系统动量中心参考系中, 一质子以 $\beta_{c} = 0.2$ , 向与前进方向的夹角 $\theta_{c} = 60^{\circ}$ 的方向发射, 求在实验室参考系中的速率 $\beta_{l} c$ 和方向 $\theta_{l}$ , 不超过几个百分点, 可作非相对论近似.
\end{problem}
\clearpage


\begin{problem}{12.4.16}
费米实验室中的高能中微子束, 是由先形成一束单能的 $\pi^{+}$ (或 $\mathrm{K}^{+}$ ) 介子束, 而后让 $\pi$ 介子衰变: $\pi^{+} \rightarrow \mu^{+} + \nu$ 而产生的, 已知 $\pi$ 介子的静止质量为 $140\mathrm{MeV} / c^2$ , $\mu$ 子的静止质量为 $106\mathrm{MeV} / c^2$ , 求:

(1) 在 $\pi^{+}$ 静止参考系中衰变得到的中微子的能量；

(2) 在实验室参考系中, 中微子的能量与衰变角 $\theta$ (见图 12.29) 有关, 若每个 $\pi^{+}$ 介子

有 200GeV 能量, 求产生在正前方( $\theta=0$ )的一个中微子的能量;

(3) 中微子能量降到最大值一半的衰变角 $\theta$ .
\begin{center}
\begin{tikzpicture}[>=Stealth, thick, line cap=round, line join=round]
  \coordinate (In) at (0, 0);
  \coordinate (O)  at (2.5, 0);
  \coordinate (F)  at (5.0, 0);
  \coordinate (Nu) at (4.7, 1.25);
  \coordinate (Mu) at (4.5, -1.0);

  % 入射介子
  \draw[postaction={decorate}, decoration={markings, mark=at position 0.52 with {\arrow{Stealth}}}]
    (In) -- (O);
  \node[below=2pt] at (0.85, 0) {$\pi^+$};

  % 原方向虚线延长线
  \draw[dashed] (O) -- (F);

  % 衰变出射粒子
  \draw[->] (O) -- (Nu) node[above left=-1pt and 1pt] {$\nu$};
  \draw[->] (O) -- (Mu) node[below=2pt] {$\mu^+$};

  % 衰变角标注
  \pic[draw, thin, angle radius=0.95cm, angle eccentricity=1.35, "$\theta$"] {angle=F--O--Nu};
  \pic[draw, thin, angle radius=0.95cm, angle eccentricity=1.35, "$\alpha$"] {angle=Mu--O--F};
\end{tikzpicture}

\end{center}
\end{problem}
\clearpage


\begin{problem}{12.4.17}
如 12.4.15 题中的靶核是固定不动的, $\beta_{0}c$ 是弹回后系统的动量中心系的速率. $E_{0}$ 是弹回后系统在动量中心参考系中的能量(设没有新粒子产生),12.4.15 的结果还能用吗?
\end{problem}
\clearpage


\begin{problem}{12.4.18}
以速度 v 运动的 $\pi^{0}$ 介子(静止质量为 m) 裂变为两个相同的 $\gamma$ 量子, 试确定 $\gamma$ 量子的飞散角.
\end{problem}
\clearpage


\begin{problem}{12.4.19}
用真空中电磁波的相位是洛伦兹变换下的不变量,导出多普勒效应的下列两个式子:

\[
\begin{array}{l} \nu = \sqrt {1 - \left(\frac {v}{c}\right) ^ {2}} \frac {c}{c - v \cos \theta} \nu^ {\prime} \\ \nu = \frac {1}{\sqrt {1 - \left(\frac {v}{c}\right) ^ {2}}} \frac {c + v \cos \theta^ {\prime}}{c} \nu^ {\prime} \end{array}
\]
\end{problem}
\clearpage


\begin{problem}{12.4.20}
以 $\beta = \frac{v}{c} = \frac{1}{\sqrt{2}}$ 运动的、每秒钟流量 $10^{6}$ 的 $\mathrm{K}_l^0$ 介子束与铅块相互作用，反应为

\[
\mathbf {K} _ {l} ^ {0} + \text { 铅块 } \rightarrow \mathbf {K} _ {s} ^ {0} + \text { 铅块 }
\]

反应前后铅块的内部状态是相同的,入射的 $K_{l}^{0}$ 与出射的 $K_{s}^{0}$ 的运动方向也是一样的.用

\[
\begin{array}{r l} & m (\mathrm{K} _ {l}) = 5 \times 1 0 ^ {8} \mathrm{eV} / c ^ {2} \\ & m (\mathrm{K} _ {l}) - m (\mathrm{K} _ {s}) = 3. 5 \times 1 0 ^ {- 6} \mathrm{eV} / c ^ {2} \end{array}
\]

求出因为这一过程作用在铅块上的平均力的大小和方向(以达因或牛顿为单位).
\end{problem}
\clearpage


\begin{problem}{12.4.21}
一个静止质量为 $m$ 、荷电 $q$ 的粒子，在一均匀磁场 $\bm{B} = B\bm{k}$ 中以半径 $R$ 在 $xy$ 平面中做圆周运动.

(1) 求 B 与 q、R、m 及角频率 $\omega$ 的关系；

(2) 因为磁场 B 对粒子不做功, 粒子的速率为常数, 但另一以匀速 Vi 运动的观察者却看到粒子的速率不是常数, 这个观察者测得的 $u_{4}^{\prime}$ [粒子四维速度的第四分量 (粒子的四维速度矢量定义为 $\gamma_{u}u, \gamma_{u}c$ )] 是什么样的?

(3) 计算 $\frac{\mathrm{d}u_4'}{\mathrm{d}\tau}, \frac{\mathrm{d}p_4'}{\mathrm{d}\tau}$ , 在观察者看来, 粒子的能量怎样改变?
\end{problem}
\clearpage


\begin{problem}{12.4.22}
一个静止质量为 $m$ 、荷电 $e$ 的粒子在恒定的电磁场中运动，在某一惯性参考系中， $E = ai, B = bk$ . 阐明粒子的四维速度作为固有时的函数的微分方程，证明方程的解是指数函数的叠加，确定这些指数，在什么条件下（关于 $\bm{E}$ 和 $\bm{B}$ ）四维速度的所有分量沿其轨道是有界的.
\end{problem}
\clearpage


\begin{problem}{12.4.23}
考虑图 12.32 所示的反应, 其中静止质量为 $m_{1}$ 和 $m_{2}$ 的粒子是入射粒子, 而静止质量为 $m_{3}$ 和 $m_{4}$ 的粒子是出射粒子, $\mathbb{k}, \mathbb{I}\mathbb{p}$ 是它们的四维动量, $\mathbb{I}\mathbb{k}$ 的定义为 $\mathbb{I}\mathbb{k} = \left(k, \mathrm{i}\frac{E}{c}\right)$ , 下面给出的变量常用来描述这种类型的反应.

\[
s = - \left(\mathbb {1} \mathrm{k} _ {1} + \mathbb {1 p} _ {1}\right) ^ {2}, t = - \left(\mathbb {1} \mathrm{k} _ {1} - \mathbb {1 k} _ {2}\right) ^ {2}, u = - \left(\mathbb {1} \mathrm{k} _ {1} - \mathbb {1 p} _ {2}\right) ^ {2}
\]

$(\mathbb{1}\mathbf{k}_1 + \mathbb{1}\mathbf{p}_1)^2$ 是 $(\mathbb{1}\mathbf{k}_1 + \mathbb{1}\mathbf{p}_1) \cdot (\mathbb{1}\mathbf{k}_1 + \mathbb{1}\mathbf{p}_1)$ 的缩写.

(1) 证明: $s + t + u = \sum_{i=1}^{4} m_i^2 c^2$ ;

(2) 假设反应是弹性散射, 并令 $m_{1} = m_{3} = \mu, m_{2} = m_{4} = m$ , 在动量中心参考系中, 令静止质量为 $\mu$ 的粒子的入射和出射的三维动量分别为 $\bm{k}$ 和 $\bm{k}'$ , 尽可能简单地用 $\bm{k}, \bm{k}'$ 将 $s, t, u$ 表示出来, 说明 $s, t, u$ 的物理意义;

(3) 假定在实验室参考系中, (2) 向中所述静止质量为 $m$ 的粒子初态是静止的, 将静止质量为 $\mu$ 的粒子在实验室参考系中的初能量、末能量及散射角用 $s, t, u$ 表示出来.
\begin{center}
\begin{tikzpicture}[>=Stealth, thick, line cap=round, line join=round,
  fermion/.style={postaction={decorate}, decoration={markings, mark=at position 0.52 with {\arrow{Stealth}}}}]
  % 中心相互作用区域（椭圆）
  \draw (0, 0) ellipse [x radius=0.9cm, y radius=0.5cm];

  % 椭圆边界连接点
  \coordinate (BL) at ({0.9*cos(225)}, {0.5*sin(225)});
  \coordinate (BR) at ({0.9*cos(-45)}, {0.5*sin(-45)});
  \coordinate (TL) at ({0.9*cos(135)}, {0.5*sin(135)});
  \coordinate (TR) at ({0.9*cos(45)},  {0.5*sin(45)});

  % 外部端点
  \coordinate (P1in)  at (-2.0, -1.45);
  \coordinate (P2in)  at (2.0, -1.45);
  \coordinate (P3out) at (-2.0, 1.45);
  \coordinate (P4out) at (2.0, 1.45);

  % 入射粒子
  \draw[fermion] (P1in) -- node[pos=0.45, above left=0pt and 2pt] {$\boldsymbol{k}_1$} (BL);
  \node[below=2pt] at (P1in) {$m_1$};

  \draw[fermion] (P2in) -- node[pos=0.45, above right=0pt and 2pt] {$\boldsymbol{p}_1$} (BR);
  \node[below=2pt] at (P2in) {$m_2$};

  % 出射粒子
  \draw[fermion] (TL) -- node[pos=0.55, above right=0pt and 2pt] {$\boldsymbol{k}_2$} (P3out);
  \node[above left=-1pt and 1pt] at (P3out) {$m_3$};

  \draw[fermion] (TR) -- node[pos=0.55, above left=0pt and 2pt] {$\boldsymbol{p}_2$} (P4out);
  \node[above right=-1pt and 1pt] at (P4out) {$m_4$};
\end{tikzpicture}

\end{center}
\end{problem}
\clearpage

